% Options for packages loaded elsewhere
\PassOptionsToPackage{unicode}{hyperref}
\PassOptionsToPackage{hyphens}{url}
\PassOptionsToPackage{dvipsnames,svgnames,x11names}{xcolor}
%
\documentclass[
  11pt,
]{article}
\usepackage{amsmath,amssymb}
\usepackage{iftex}
\ifPDFTeX
  \usepackage[T1]{fontenc}
  \usepackage[utf8]{inputenc}
  \usepackage{textcomp} % provide euro and other symbols
\else % if luatex or xetex
  \usepackage{unicode-math} % this also loads fontspec
  \defaultfontfeatures{Scale=MatchLowercase}
  \defaultfontfeatures[\rmfamily]{Ligatures=TeX,Scale=1}
\fi
\usepackage{lmodern}
\ifPDFTeX\else
  % xetex/luatex font selection
  \setmainfont[]{Liberation Serif}
  \setsansfont[]{Liberation Sans}
  \setmonofont[]{DejaVu Sans Mono}
\fi
% Use upquote if available, for straight quotes in verbatim environments
\IfFileExists{upquote.sty}{\usepackage{upquote}}{}
\IfFileExists{microtype.sty}{% use microtype if available
  \usepackage[]{microtype}
  \UseMicrotypeSet[protrusion]{basicmath} % disable protrusion for tt fonts
}{}
\makeatletter
\@ifundefined{KOMAClassName}{% if non-KOMA class
  \IfFileExists{parskip.sty}{%
    \usepackage{parskip}
  }{% else
    \setlength{\parindent}{0pt}
    \setlength{\parskip}{6pt plus 2pt minus 1pt}}
}{% if KOMA class
  \KOMAoptions{parskip=half}}
\makeatother
\usepackage{xcolor}
\usepackage[margin=0.8in]{geometry}
\usepackage{longtable,booktabs,array}
\usepackage{calc} % for calculating minipage widths
% Correct order of tables after \paragraph or \subparagraph
\usepackage{etoolbox}
\makeatletter
\patchcmd\longtable{\par}{\if@noskipsec\mbox{}\fi\par}{}{}
\makeatother
% Allow footnotes in longtable head/foot
\IfFileExists{footnotehyper.sty}{\usepackage{footnotehyper}}{\usepackage{footnote}}
\makesavenoteenv{longtable}
\setlength{\emergencystretch}{3em} % prevent overfull lines
\providecommand{\tightlist}{%
  \setlength{\itemsep}{0pt}\setlength{\parskip}{0pt}}
\setcounter{secnumdepth}{-\maxdimen} % remove section numbering
\ifLuaTeX
\usepackage[bidi=basic]{babel}
\else
\usepackage[bidi=default]{babel}
\fi
\babelprovide[main,import]{american}
\ifPDFTeX
\else
\babelfont{rm}[]{Liberation Serif}
\fi
% get rid of language-specific shorthands (see #6817):
\let\LanguageShortHands\languageshorthands
\def\languageshorthands#1{}
\usepackage{adjustbox}
\emergencystretch=3em
\AtBeginEnvironment{longtable}{\small\sloppy}
\usepackage{stmaryrd}
% Printing shims for run-on conditional-bar tokens in the supplied Biology source.
\providecommand{\midL}{\mid L}
\providecommand{\midN}{\mid N}
\sloppy
\ifLuaTeX
  \usepackage{selnolig}  % disable illegal ligatures
\fi
\usepackage{bookmark}
\IfFileExists{xurl.sty}{\usepackage{xurl}}{} % add URL line breaks if available
\urlstyle{same}
\hypersetup{
  pdflang={en-US},
  colorlinks=true,
  linkcolor={blue},
  filecolor={Maroon},
  citecolor={Blue},
  urlcolor={blue},
  pdfcreator={LaTeX via pandoc}}

\author{}
\date{}

\begin{document}

\section{COHERENCE UNIVERSALISM:
ETHICS}\label{coherence-universalism-ethics}

\subsection{Truth, Self-Deception, Power, and
Repair}\label{truth-self-deception-power-and-repair}

\textbf{Gaura Kiśora Dās Rader}\\
\textbf{Heaven≡Earth Foundation}\\
\textbf{White Paper Series}\\
\textbf{Version 0.10}\\
\textbf{October 2026}

\textbf{Release status.} This manuscript is part of the complete Version
0.10 proposal. Numbered principles identify the framework's current
formulations; their claim status, external validation, and authorization
for consequential application remain separate questions.

\textbf{Freeze status.} This manuscript is synchronized to the frozen
Version 0.10 release registry, evidence package, and review gates.
Future changes require a new versioned release.

\section{Abstract}\label{abstract}

Ethics begins where epistemology becomes personal. The question is no
longer only whether our descriptions of reality are true. It is whether
we are willing and able to know the truth about what we are doing, why
we are doing it, whom it affects, what we benefit from, and what repair
requires. Coherence Universalism locates this question within the
Coherence Mandala: once reality includes beings with organism-relative
stakes, experience, agency, vulnerability, dependence, and unequal power
over one another, ethical relations become constitutive features of
higher-order life. The Ethics paper therefore does not attempt to derive
morality from physics or prove moral reality into existence. It asks why
agents and institutions repeatedly lose truthful contact with a moral
reality they already inhabit.

The central danger is self-deception: motivated asymmetric information
processing through which identity, attachment, loyalty, advantage,
authority, or policy influences which facts, reasons, alternatives, and
affected parties are admitted into judgment. This paper retains
justificatory integrity as indispensable while rejecting the claim that
integrity alone is sufficient for permissibility. A transparent tyrant
can accurately state a harmful objective. A sincere official can
consistently enforce an unjust rule. An institution can acknowledge the
burdens it imposes while denying affected people standing, remedy, or a
meaningful alternative. Ethical judgment therefore requires both
truthfulness in justification and substantive attention to standing,
agency, welfare, consent, rights, distribution, power, responsibility,
uncertainty, reversibility, and repair.

The paper develops a non-scalar ethical architecture. A declared problem
identifies affected beings, possible trajectories, the options actually
presented, causal assumptions, protected conditions, power relations,
benefit and burden distributions, the justificatory record, historical
path dependence, horizons, observation channels, and practical
capacities for correction. Evaluation proceeds through a profile rather
than a universal moral score. Protected floors constrain comparison;
scalarization remains conditional on disclosed weights, normalizations,
affected scope, and authority. Passing a screen establishes only
screening-admissibility relative to the represented option set, not the
adequacy of that set or all-things-considered moral permission.

The analysis then follows the truth-centered arc of the series. It
examines the anatomy of self-serving justification, explains why
introspection is insufficient, and shows why truth requires
affected-party testimony, independent records, protected dissent,
counterpower, option-set review, and institutions capable of practical
correction. It distinguishes patiency, agency, responsibility, and
sovereignty; separates causal contribution from blame and remedial duty;
treats moral distress, moral injury, and institutional betrayal without
collapsing them into one diagnosis; and reconstructs repair through
protection, truth, accountability, restitution, restoration of standing,
and prevention of recurrence. Artificial intelligence is treated as a
possible aid to moral inquiry and a source of institutional power, not
as a sovereign moral authority.

The paper culminates in \textbf{CU-COM12 --- Ethical Truthfulness and
Accountability}: the commitment to seek truth about actions before
seeking innocence for agents or institutions. Ethics completes the
corpus's return to truth. Epistemology teaches us not to lie about the
universe. Ethics teaches us not to lie about ourselves.

\textbf{Keywords:} ethics; truth; self-deception; justificatory
integrity; moral standing; power; responsibility; moral injury;
institutional betrayal; repair; artificial intelligence; Coherence
Universalism

\subsection{Relation to the series}\label{relation-to-the-series}

This paper presupposes the framework definitions and register discipline
of \emph{Coherence Universalism: Metaphysics and Epistemology} and
\emph{Foundations}, the distinction between functional stake and
phenomenal valence developed across Biology and Consciousness, the
account of meaning-bearing regulation and self-deception developed in
Psychology, and the analysis of institutions, power, dissent, and
practical correctability developed in Social Dynamics. It is written to
stand on its own, but its full architecture is relational: the
scientific and philosophical papers identify the kinds of beings and
dependencies that make ethical relations possible, while this paper asks
what truthfulness, protection, accountability, and repair require once
those relations exist.

\subsection{Note on Version 0.10}\label{note-on-version-0.10}

Version 0.10 preserves the original paper's insight that moral reasoning
is vulnerable to motivated corruption while replacing the earlier
single-principle architecture with a plural constraint framework. Full
revision history, formal contrast models, claim ledgers, independent
audit materials, and reproducibility files are preserved in the
\href{../supplements/08_Ethics_Technical_Supplement_v0.10.md}{technical
supplement} and archive.

\textbf{Release-candidate synchronization.} This manuscript uses the
shared Version 0.10 definitions, principle identifiers, Mandala
cross-references, and interpretive and commitment language published in
\emph{Coherence Universalism: Metaphysics and Epistemology} and
\emph{Coherence Universalism: Foundations}. Its audited technical core,
formal results, numerical results, and claim-status boundaries remain
unchanged except where this manuscript explicitly states otherwise.

\subsection{How claims are marked}\label{how-claims-are-marked}

Definitions, propositions, models, hypotheses, interpretive claims, and
commitment-register statements are identified where their status
matters. Normative principles are not presented as empirical
discoveries. Empirical findings constrain application; they do not by
themselves determine the content of rights, protected floors,
legitimacy, or repair.

\section{Contents}\label{contents}

\begin{enumerate}
\def\labelenumi{\arabic{enumi}.}
\tightlist
\item
  Ethics begins where epistemology becomes personal\\
\item
  Declaring an ethical problem\\
\item
  Why moral agents lose contact with truth\\
\item
  Justificatory integrity: necessary, not sufficient\\
\item
  Truth requires more than introspection\\
\item
  Compassion as an epistemic discipline\\
\item
  Standing, patiency, agency, and rights\\
\item
  Consequences, distribution, uncertainty, and tragic conflict\\
\item
  Power, coercion, legitimacy, and public justification\\
\item
  Responsibility under role and distributed causation\\
\item
  Moral distress, moral injury, and institutional betrayal\\
\item
  Accountability, repair, and non-recurrence\\
\item
  Artificial intelligence, moral assistance, and sovereignty\\
\item
  Research program and principle architecture\\
\item
  What commitments this account licenses\\
\item
  Conclusion: truth before innocence\\
  Appendix A. Ethics principles and application governance\\
  Appendix B. Major revisions from v0.09\\
  Appendix C. Technical supplement and reproducibility
\end{enumerate}

\section{Principal Contributions}\label{principal-contributions}

\begin{enumerate}
\def\labelenumi{\arabic{enumi}.}
\tightlist
\item
  \textbf{Truth-centered ethics.} Ethics is organized around truthful
  contact with motive, consequence, power, benefit, omission, and repair
  rather than around a new proof that moral claims exist.
\item
  \textbf{Integrity without sufficiency.} Justificatory integrity is
  necessary for defensible moral judgment but cannot by itself establish
  permissibility.
\item
  \textbf{Plural ethical profile.} Standing, agency, welfare, consent,
  distribution, legitimacy, responsibility, reversibility, coercion, and
  repair remain distinct rather than being collapsed into one scalar.
\item
  \textbf{Option-set integrity.} Ethical screening and ranking are
  conditional on the alternatives permitted to appear.
\item
  \textbf{Truth beyond introspection.} Affected-party testimony,
  independent records, dissent, counterpower, and practical
  correctability are epistemic requirements where power is unequal.
\item
  \textbf{Distributed responsibility.} Causal contribution, knowledge,
  control, authority, benefit, blame, liability, and remedial duty are
  evaluated separately.
\item
  \textbf{Repair after truth.} Acknowledgment is joined to protection,
  record correction, restitution, restored standing, structural change,
  and non-recurrence.
\item
  \textbf{Patiency without sovereignty.} Welfare-relevant experience can
  ground protection without conferring responsibility or legitimate
  authority.
\item
  \textbf{AI as moral assistant, not sovereign.} AI may expose
  inconsistency and omitted perspectives while remaining embedded in
  institutional power and human responsibility.
\item
  \textbf{Commitment to ethical truthfulness.} CU-COM12 makes truth
  before innocence an official commitment of the framework.
\end{enumerate}

\begin{center}\rule{0.5\linewidth}{0.5pt}\end{center}

\section{Domain Analytic Frame and CAP
Record}\label{domain-analytic-frame-and-cap-record}

\textbf{Target:} Coherence Universalism Version 0.10 release\\
\textbf{Status:} Release declaration; scope-limited and revisable under
the declared review gates\\
\textbf{Use rule:} Unresolved fields are left explicit rather than
completed by analogy.

\subsection{Paper and model identity}\label{paper-and-model-identity}

\begin{longtable}[]{@{}
  >{\raggedright\arraybackslash}p{(\columnwidth - 4\tabcolsep) * \real{0.3333}}
  >{\raggedright\arraybackslash}p{(\columnwidth - 4\tabcolsep) * \real{0.3333}}
  >{\raggedright\arraybackslash}p{(\columnwidth - 4\tabcolsep) * \real{0.3333}}@{}}
\toprule\noalign{}
\begin{minipage}[b]{\linewidth}\raggedright
Field
\end{minipage} & \begin{minipage}[b]{\linewidth}\raggedright
Declaration
\end{minipage} & \begin{minipage}[b]{\linewidth}\raggedright
Status
\end{minipage} \\
\midrule\noalign{}
\endhead
\bottomrule\noalign{}
\endlastfoot
Paper & Coherence Universalism: Ethics & Specified \\
Domain owner & Ethics & Specified by the Version 0.10 release corpus \\
Primary model or analysis & Ethical problem
\(\mathcal E=(A,\Omega,U,F,K,P,R,J,\mathcal H,\pi,T,\mathcal O,\mathcal C)\)
& Paper-relative \\
Primary claim status & Normative principles, definitions, empirical
imports, bridge principles, diagnostic hypotheses, and commitments &
Mixed statuses remain explicit \\
Review gate & Ethics, law, political philosophy, social science,
relevant domain expertise, and heightened affected-community
co-governance & Proportionate to claim and consequential reach \\
Methodological authority & CU-FP2, CU-M2, CU-M11, CU-M12, CU-M13,
CU-M14, and CU-M15 & Controlling Version 0.10 methodological
framework \\
\end{longtable}

\subsection{Shared Analytic Frame}\label{shared-analytic-frame}

\begin{longtable}[]{@{}
  >{\raggedright\arraybackslash}p{(\columnwidth - 4\tabcolsep) * \real{0.3333}}
  >{\raggedright\arraybackslash}p{(\columnwidth - 4\tabcolsep) * \real{0.3333}}
  >{\raggedright\arraybackslash}p{(\columnwidth - 4\tabcolsep) * \real{0.3333}}@{}}
\toprule\noalign{}
\begin{minipage}[b]{\linewidth}\raggedright
\#
\end{minipage} & \begin{minipage}[b]{\linewidth}\raggedright
Slot and domain declaration
\end{minipage} & \begin{minipage}[b]{\linewidth}\raggedright
Status and evidence
\end{minipage} \\
\midrule\noalign{}
\endhead
\bottomrule\noalign{}
\endlastfoot
1 & \textbf{Unit and boundary.} A declared ethical problem containing
affected beings, trajectories, presented and feasible options,
institutions, histories, and jurisdictions. & Affected-being scope is
load-bearing \\
2 & \textbf{State or trajectory space.} \(\Omega\) contains possible
states or trajectories; ethical evaluation often concerns action paths
and repair, not static outcomes alone. & Declared \\
3 & \textbf{Environment.} Material, social, legal, ecological,
institutional, and historical conditions shape feasibility and
consequence. & Included through the causal model \\
4 & \textbf{Constraints and viable region.} \(K\) contains rights,
floors, prohibitions, standing, and emergency-exception rules; these are
not high weights. & Screening architecture \\
5 & \textbf{Actions, options, and transitions.}
\(U^{pres}\subseteq U^{feas}\) distinguishes presented from feasible
acts, omissions, policies, delay, redesign, and repair. & Option-set
integrity required \\
6 & \textbf{Dynamics or update rule.} \(F\) is a descriptive causal
model linking actions to consequences; Ethics asserts no natural moral
flow or universal optimizer. & Causal and trajectory-specific \\
7 & \textbf{Observation and access.} \(\mathcal O\), justificatory
record \(J\), testimony, independent evidence, and uncertainty determine
what enters judgment. & Truth and provenance required \\
8 & \textbf{Perturbation and intervention class.} Policy change,
refusal, emergency action, review, rollback, restitution, restoration,
and prevention of recurrence are distinct interventions. &
Consequentially governed \\
9 & \textbf{Equilibria, basins, or regimes.} No natural ethical
equilibrium or basin is claimed. Stable institutions can remain unjust.
& Not applicable generally \\
10 & \textbf{Local--global status.} Aggregate and local effects are
separated; normative noncompensation comes from protected standing and
floors, not from convexity or information geometry. & Qualitative /
coordinate-wise unless separately proved \\
11 & \textbf{Scalar status.} The primary object is a plural ethical
profile. Post-screen decision scalarization is permitted only with
disclosed authority and tradeoffs. & No universal moral score \\
12 & \textbf{Statistical-image and allocation status.} Outcome, burden,
risk, and resource distributions may qualify. Standing, rights,
justification, responsibility, and the ethical profile are not simplex
coordinates. & Statistical subobjects only \\
13 & \textbf{Translation status.} Empirical and formal methods inform
but do not determine rights, standing, legitimacy, or moral obligation
without explicit normative bridges. & Bounded \\
14 & \textbf{Open gaps.} Option-set adequacy, authority to set floors,
moral uncertainty, standing across unfamiliar beings, responsibility
under distributed causation, and repair remain contested. & Open
normative work \\
\end{longtable}

\subsection{Compact Coherence Analysis Protocol
record}\label{compact-coherence-analysis-protocol-record}

\begin{longtable}[]{@{}
  >{\raggedright\arraybackslash}p{(\columnwidth - 4\tabcolsep) * \real{0.3333}}
  >{\raggedright\arraybackslash}p{(\columnwidth - 4\tabcolsep) * \real{0.3333}}
  >{\raggedright\arraybackslash}p{(\columnwidth - 4\tabcolsep) * \real{0.3333}}@{}}
\toprule\noalign{}
\begin{minipage}[b]{\linewidth}\raggedright
CAP module
\end{minipage} & \begin{minipage}[b]{\linewidth}\raggedright
Declaration
\end{minipage} & \begin{minipage}[b]{\linewidth}\raggedright
Boundary
\end{minipage} \\
\midrule\noalign{}
\endhead
\bottomrule\noalign{}
\endlastfoot
System declaration & Name affected beings, trajectories, presented and
feasible options, causal model, power, history, boundary, and horizon. &
Section 2.1 \\
Observation declaration & Disclose evidence, testimony, missing parties,
uncertainty, provenance, and who controls the record. & Self-report and
institutional records are corrigible \\
Evaluation declaration & Screen floors first; then compare admissible
options with profiles, distributions, uncertainty, and preserved
incomparability. & Passing a screen is not moral permission \\
Power and authorship declaration & Name who framed the question,
generated options, classified feasibility, set floors or weights,
benefits, bears harm, and controls appeal or remedy. & Central ethical
requirement \\
Revision declaration & New facts within fixed values differ from changed
standing, floors, option set, horizon, or authority. & Changed model
requires renewed justification \\
\end{longtable}

\subsection{License summary}\label{license-summary}

\begin{longtable}[]{@{}
  >{\raggedright\arraybackslash}p{(\columnwidth - 4\tabcolsep) * \real{0.3333}}
  >{\raggedright\arraybackslash}p{(\columnwidth - 4\tabcolsep) * \real{0.3333}}
  >{\raggedright\arraybackslash}p{(\columnwidth - 4\tabcolsep) * \real{0.3333}}@{}}
\toprule\noalign{}
\begin{minipage}[b]{\linewidth}\raggedright
License
\end{minipage} & \begin{minipage}[b]{\linewidth}\raggedright
Declaration
\end{minipage} & \begin{minipage}[b]{\linewidth}\raggedright
Current status
\end{minipage} \\
\midrule\noalign{}
\endhead
\bottomrule\noalign{}
\endlastfoot
Scalar license & Only purpose-bound post-screen decision scalarizations
or construct measures; no universal moral scalar. & Conditional \\
Local--global license & Normative noncompensation plus qualitative or
coordinate-wise distribution analysis; exact math only for licensed
subobjects. & Mixed \\
Statistical Image License & Possible for uncertainty, outcome, burden,
exposure, or risk distributions. & Object-specific \\
Allocation License & Possible for common-unit resources or burdens;
protected standing and floors remain outside compensatory allocation
arithmetic. & Conditional \\
Refinement and Coarse-Graining License & Aggregation is licensed only
when affected beings and morally relevant distinctions are not erased. &
Heightened review \\
Fisher-Rao License & Possible for statistical subobjects only; geometry
supplies no moral direction, standing, or legitimacy. & No normative
inference \\
Translation license & Descriptive findings can constrain but do not
derive normative conclusions. & Explicit bridge required \\
\end{longtable}

\subsection{Open gaps and review
disposition}\label{open-gaps-and-review-disposition}

\begin{longtable}[]{@{}
  >{\raggedright\arraybackslash}p{(\columnwidth - 4\tabcolsep) * \real{0.3333}}
  >{\raggedright\arraybackslash}p{(\columnwidth - 4\tabcolsep) * \real{0.3333}}
  >{\raggedright\arraybackslash}p{(\columnwidth - 4\tabcolsep) * \real{0.3333}}@{}}
\toprule\noalign{}
\begin{minipage}[b]{\linewidth}\raggedright
Gap
\end{minipage} & \begin{minipage}[b]{\linewidth}\raggedright
Consequence
\end{minipage} & \begin{minipage}[b]{\linewidth}\raggedright
Required work
\end{minipage} \\
\midrule\noalign{}
\endhead
\bottomrule\noalign{}
\endlastfoot
Option adequacy & A correct procedure can choose from the wrong menu. &
Mechanisms for omitted-option discovery and affected-party proposals \\
Floor authority & Thresholds can encode domination or exclusion. &
Public justification and affected-community authority \\
Repair evidence & Apology or record correction may not restore standing
or prevent recurrence. & Track protection, remedy, restitution, and
recurrence \\
\end{longtable}

\begin{quote}
\textbf{Propagation rule.} This card makes the paper's existing objects
and limits visible. It does not strengthen any empirical, mathematical,
normative, metaphysical, or spiritual claim beyond the wording supported
in the manuscript.
\end{quote}

\begin{center}\rule{0.5\linewidth}{0.5pt}\end{center}

\section{1. Ethics begins where epistemology becomes
personal}\label{ethics-begins-where-epistemology-becomes-personal}

\subsection{1.1 Truth about reality and truth about
participation}\label{truth-about-reality-and-truth-about-participation}

Epistemology asks what we are entitled to believe and say about reality.
Ethics asks what prevents us from seeing truthfully what we are doing
within that reality.

The truth at stake is not only the accuracy of a proposition. It
includes the truth about motive, dependence, benefit, power, omission,
consequence, and responsibility. It asks whether an agent who condemns
exploitation can acknowledge the labor from which they benefit; whether
an institution that celebrates transparency preserves the records by
which its conduct could be challenged; whether a movement that speaks in
the name of justice applies the same evidential standards to its allies
and opponents; and whether a person who has caused harm is seeking to
understand what happened or merely to preserve an identity of innocence.

Ethics therefore begins where epistemology becomes personal. A false
factual premise can corrupt moral judgment, but moral distortion often
persists even when the relevant facts are available. What fails is the
structure through which facts, reasons, testimony, and alternatives are
admitted. Identity, loyalty, shame, fear, status, attachment,
resentment, and advantage can determine what becomes visible and what is
neutralized before it reaches judgment.

This paper calls that family of processes \textbf{self-deception} when
the asymmetry is motivated and recurrent. The term does not imply
conscious lying. It names a way in which persons and institutions can
remain sincere while selectively protecting themselves from reasons that
would require revision.

\subsection{1.2 Moral reality inherited from the
Mandala}\label{moral-reality-inherited-from-the-mandala}

Coherence Universalism does not begin Ethics by attempting to derive
obligation from a physical equation. Physics describes lawful
possibility. Biology identifies active self-maintenance and
organism-relative stake. Consciousness introduces the possibility of
felt experience. Psychology develops meaning, identity, motivation, and
self-deception. Social Dynamics examines coordination, power,
institutions, dissent, and repair.

The Coherence Mandala interprets these domains as reciprocally dependent
aspects of one reality. Once that reality contains beings for whom
conditions can be better or worse, centers of experience capable of
suffering, agents able to act for reasons, relationships of dependency,
and institutions with unequal power over lives, ethical relations are
not optional decorations on an otherwise complete world. They are
constitutive features of higher-order coexistence.

This is a systems-metaphysical and normative bridge, not a deduction of
morality from physics. It supplies the ground from which the Ethics
paper proceeds. The central problem is not whether moral reality exists,
but whether agents and institutions remain truthful enough to
participate in it responsibly.

In the canonical sixty-four-stage Coherence Mandala, the central ethical
sequence occupies Stages 27--31: Value Differentiation, Harm
Recognition, Tradeoffs Under Protected Floors, Justification and Public
Reason, and Ethical Institutions and Repair. The later reflexive
stages---especially Universal Moral Concern, Coherence Stewardship,
Open-Ended Stability, Meaning Without Finality, and Stage 64, Coherence
Universalism---show why ethical recognition culminates in responsibility
rather than in theoretical closure.

\subsection{1.3 Why self-deception is
central}\label{why-self-deception-is-central}

Most serious moral failures are not committed by agents who possess no
moral language. They are committed through moral language: security,
efficiency, duty, loyalty, progress, necessity, care, order, freedom,
responsibility. The same vocabulary can name a genuine reason or conceal
a transfer of cost.

Self-deception matters because it corrupts the conditions under which
reasons can function. It can:

\begin{itemize}
\tightlist
\item
  narrow the affected-party set until the people bearing the burden
  disappear;
\item
  redefine coercion as consent because resistance is absent;
\item
  treat presented options as exhaustive because alternatives threaten
  authority;
\item
  change evidential standards when a conclusion protects identity or
  advantage;
\item
  convert structural harm into individual failure;
\item
  redescribe punishment as repair;
\item
  and use uncertainty selectively---demanding impossible proof from the
  harmed while accepting weak reassurance from the powerful.
\end{itemize}

The wrong is not merely that a conclusion may be false. The process
disables correction while preserving the appearance of justification.

\subsection{1.4 Justificatory integrity is necessary but not
sufficient}\label{justificatory-integrity-is-necessary-but-not-sufficient}

A defense that depends upon distorted facts, hidden alternatives,
excluded testimony, or asymmetric standards is morally defective. This
is the \textbf{Justificatory Integrity Constraint}.

But integrity alone is not enough. An actor can openly state that
another group will be used as a means. A policy can be transparent,
internally consistent, and devastating. A system can record every burden
it imposes while denying affected people standing to challenge the
objective itself.

Ethical judgment must therefore face substantive questions that
truthfulness cannot answer by itself:

\begin{itemize}
\tightlist
\item
  Who possesses moral standing?
\item
  Which interests or capabilities deserve protected status?
\item
  What kinds of consent are valid under dependency?
\item
  Which harms may not be compensated by unrelated gains?
\item
  How should uncertainty and irreversibility alter action?
\item
  Who has legitimate authority?
\item
  How should responsibility be distributed?
\item
  What does repair require after wrongdoing?
\end{itemize}

CU Ethics is consequently a plural constraint framework, not a theory
that reduces morality to sincerity or coherence.

\subsection{1.5 Moral pluralism without moral
quietism}\label{moral-pluralism-without-moral-quietism}

The architecture draws from several traditions without claiming that
they say the same thing. Deontological ethics protects persons from
being treated merely as instruments. Consequentialism insists that
outcomes matter. Contractualism and public-reason traditions demand
answerability to those governed. Care ethics makes dependency,
embodiment, and relationship morally visible. Capability and rights
approaches protect substantive possibilities rather than formal
permission alone. Virtue ethics asks what kinds of agents and practices
make good judgment possible. Feminist, disability, decolonial, and
critical traditions reveal how apparently neutral abstractions can
encode dominant standpoints. Restorative and transitional justice
distinguish repair from punishment and private remorse.

CU coordinates these insights around declared affected parties,
protected conditions, power, responsibility, truthfulness, and repair.
It does not settle every disagreement among them. Where moral reasons
remain in genuine conflict, the framework should display the conflict
rather than manufacture a scalar answer.

\subsection{1.6 Scope and roadmap}\label{scope-and-roadmap}

The paper offers:

\begin{itemize}
\tightlist
\item
  a declaration discipline for ethical problems;
\item
  a multidimensional profile and protected-floor architecture;
\item
  a truth-centered account of self-deception and justificatory
  integrity;
\item
  substantive treatments of standing, consequences, power,
  responsibility, moral injury, and repair;
\item
  a bounded account of contemplative moral practice;
\item
  and an analysis of artificial intelligence as both moral assistance
  and institutional power.
\end{itemize}

It does not provide a universal algorithm for right action, a substitute
for law or democratic judgment, or a license to diagnose opponents. Its
principal aim is to make moral reasoning more truthful, more accountable
to affected beings, and more capable of repair.

\section{2. Declaring an ethical
problem}\label{declaring-an-ethical-problem}

\subsection{2.1 The ethical model tuple}\label{the-ethical-model-tuple}

A candidate ethical problem is represented by

\[\adjustbox{max width=\linewidth}{$\displaystyle 
\mathcal E
=
(A,\Omega,U,F,K,P,R,J,\mathcal H,\pi,T,\mathcal O,\mathcal C).
$}\]

The components are:

\begin{itemize}
\tightlist
\item
  \(A\): affected beings, groups, representatives, and future or
  indirect claimants;
\item
  \(\Omega\): possible states, trajectories, or scenarios;
\item
  \(U^{\mathrm{pres}}\): the actions, omissions, policies, and
  institutional responses actually presented for comparison, understood
  as a possibly incomplete subset of the feasible alternatives
  \(U^{\mathrm{feas}}\);
\item
  \(F\): descriptive and causal model connecting actions to
  consequences;
\item
  \(K\): protected conditions, rights, floors, prohibitions, and
  emergency exceptions;
\item
  \(P\): power, authority, dependency, coercion, ownership, and exit
  structure;
\item
  \(R\): distributions of benefit, burden, risk, voice, and repair
  obligation;
\item
  \(J\): justificatory record, including reasons, evidence,
  alternatives, uncertainty, provenance, and dissent;
\item
  \(\mathcal H\): history, prior commitments, inherited harms, and path
  dependence;
\item
  \(\pi\): boundary, jurisdiction, scale, and representation map;
\item
  \(T\): relevant temporal horizons;
\item
  \(\mathcal O\): observation, measurement, and reporting channels; and
\item
  \(\mathcal C\): correction, appeal, rollback, remedy, and repair
  capacities.
\end{itemize}

The tuple is not a claim that moral life can be fully formalized. It is
a declaration discipline. It prevents an analyst from evaluating ``the
policy'' while leaving unnamed the people excluded from the model, the
institution that controls the evidence, or the losses that occur after
the official time horizon ends.

The standardized Domain Analytic Frame above is the compact CAP record
for this paper. It preserves the distinction between mathematical
relation status under proposed CU-M13 and normative noncompensation
under CU-G1 and the ethical architecture.

The distinction between feasible and presented options is load-bearing:

\[\adjustbox{max width=\linewidth}{$\displaystyle 
U^{\mathrm{pres}}\subseteq U^{\mathrm{feas}}.
$}\]

A decision procedure can compare only the alternatives represented in
\(U^{\mathrm{pres}}\). It cannot infer from the success of one
represented option that the set is adequate, that a less restrictive or
more reparative alternative is unavailable, or that the institution
controlling option generation has searched in good faith.

\subsection{2.2 The affected-being set}\label{the-affected-being-set}

The set \(A\) is normatively load-bearing. It must not be limited
automatically to:

\begin{itemize}
\tightlist
\item
  formal decision makers;
\item
  consenting participants;
\item
  citizens of one jurisdiction;
\item
  present persons;
\item
  property owners;
\item
  humans capable of verbal reason;
\item
  or beings whose harms are easy to measure.
\end{itemize}

Potentially affected parties include direct subjects, families, workers,
nonusers, data subjects, displaced communities, future populations,
nonhuman animals, ecosystems as legally or morally represented, and
possible artificial patients where credible uncertainty exists.

Inclusion does not determine equal weight or identical rights. It
prevents absence from being mistaken for irrelevance.

\subsection{2.3 A plural ethical
profile}\label{a-plural-ethical-profile}

the analysis uses the profile

\[\adjustbox{max width=\linewidth}{$\displaystyle 
\mathbf P_E(u)
=
\langle
S,A_g,W,C_n,D,J_i,L,p_{\mathrm{resp}},U_r,Q,R_p
\rangle_u.
$}\]

The coordinates are candidate dimensions rather than universally fixed
measures:

\begin{itemize}
\tightlist
\item
  \(S\): standing and recognition of affected beings;
\item
  \(A_g\): protection of agency, choice, and substantive capability;
\item
  \(W\): welfare, suffering, need, and morally relevant interest;
\item
  \(C_n\): consent, authorization, and freedom from undue influence;
\item
  \(D\): distribution of benefit, burden, risk, and sacrifice;
\item
  \(J_i\): justificatory integrity and evidential truthfulness;
\item
  \(L\): legitimacy and public answerability;
\item
  \(R_a\): alignment of responsibility with knowledge, control,
  authority, and benefit;
\item
  \(U_r\): uncertainty management, precaution, reversibility, and
  learning;
\item
  \(Q\): coercion, domination, extraction, and dependency costs; and
\item
  \(R_p\): protection, correctability, remedy, and repair.
\end{itemize}

A profile may support threshold screening, Pareto comparison,
lexicographic priority, scenario analysis, or a declared scalarization.
It does not automatically support a total ranking.

\subsection{2.4 Protected floors and secondary
comparison}\label{protected-floors-and-secondary-comparison}

Let \(k_j\) be a protected floor for dimension \(j\), and let
\(y_{aj}(u,\omega)\) be the outcome for affected being or group \(a\) in
scenario \(\omega\). A simple screening rule is

\[\adjustbox{max width=\linewidth}{$\displaystyle 
u\in U_K
\quad\Longleftrightarrow\quad
\Pr\!\left[
 y_{aj}(u,\omega)\ge k_{aj}
 \text{ for all protected }(a,j)
\right]
\ge 1-\delta,
$}\]

subject to the chosen uncertainty model.

This is only a formal shell. Ethics must still justify:

\begin{itemize}
\tightlist
\item
  which dimensions receive floor status;
\item
  whose floors count;
\item
  the threshold values;
\item
  the acceptable probability of violation;
\item
  whether worst-case, expected, or distributionally robust criteria are
  appropriate;
\item
  and what emergency exceptions exist.
\end{itemize}

After screening, options may remain partially ordered. A secondary
decision rule may then be used, but it must not silently compensate for
what the floor was meant to protect.

\subsubsection{2.4.1 Conditional admissibility and option-set
adequacy}\label{conditional-admissibility-and-option-set-adequacy}

Let

\[\adjustbox{max width=\linewidth}{$\displaystyle 
U_K^{\mathrm{pres}}
=
U^{\mathrm{pres}}\cap U_K.
$}\]

Then \(u\in U_K^{\mathrm{pres}}\) means only that \(u\) passes the
declared screen relative to the represented parties, evidence,
scenarios, floors, and option set. this account calls this
\textbf{screening-admissible relative to the presented set}. It is not
an all-things-considered declaration of moral permissibility.

An option-set adequacy review should ask, where relevant:

\begin{itemize}
\tightlist
\item
  who generated and excluded alternatives;
\item
  whether affected parties could introduce options;
\item
  whether the status quo, no-action, delay, staged action,
  least-restrictive action, and repair or redesign options were
  considered;
\item
  whether legal, financial, temporal, or institutional constraints were
  treated as fixed when they could themselves be changed;
\item
  whether the apparent choice was produced by prior coercion, capture,
  or path dependence; and
\item
  what evidence supports the claim that omitted alternatives are
  infeasible.
\end{itemize}

If every presented option fails, the result may be tragedy, emergency
minimization, an inadequate model, or a failed option-generation
process. If exactly one presented option passes, it is uniquely
screening-admissible \textbf{within that set}, not thereby morally
endorsed against alternatives that were never allowed to appear.

\subsubsection{2.4.2 Reversibility under credible severe-harm
risk}\label{reversibility-under-credible-severe-harm-risk}

Reversibility is not a universal trump. Surgery, rescue, testimony,
parenthood, ecological intervention, and delayed action can each contain
irreversible consequences. Yet when an application stipulates a credible
risk of severe or irreversible harm, silence about reversibility is not
neutral. The evaluator must either:

\begin{enumerate}
\def\labelenumi{\arabic{enumi}.}
\tightlist
\item
  declare a context-appropriate protected floor for reversibility,
  containment, rollback, staged exposure, or equivalent retained option
  value; or
\item
  state an exception argument showing why necessity, proportionality,
  lack of safer alternatives, and the harms of delay justify proceeding
  despite limited reversibility.
\end{enumerate}

The floor or exception remains a normative commitment. Mathematics can
apply it consistently but cannot determine its content.

\subsection{2.5 Noncompensation}\label{noncompensation}

\textbf{{[}Normative principle{]} Noncompensation} means that some
losses are not made permissible merely by adding unrelated benefits
elsewhere.

This does not imply that no burdens can ever be imposed. It requires
that burdens be examined for:

\begin{itemize}
\tightlist
\item
  severity;
\item
  necessity;
\item
  consent;
\item
  distribution;
\item
  reversibility;
\item
  alternatives;
\item
  compensation;
\item
  legitimacy;
\item
  and whether the burden destroys the standing needed to contest it.
\end{itemize}

The principle is especially important where the beneficiaries control
the measurement of aggregate benefit.

\subsection{2.6 Ethical scalarization}\label{ethical-scalarization}

A declared scalarization may take the form

\[\adjustbox{max width=\linewidth}{$\displaystyle 
V_w(u)=\sum_j w_j z_j(u),
\qquad
w_j\ge0,
\qquad
\sum_jw_j=1,
$}\]

where the \(z_j\) are normalized dimensions. Such a score can be useful
for a specific decision after protected constraints are applied. It does
not become a universal value function.

The analyst must disclose:

\begin{itemize}
\tightlist
\item
  normalization;
\item
  weights;
\item
  uncertainty treatment;
\item
  affected-being aggregation;
\item
  exclusions;
\item
  protected floors;
\item
  and the authority that chose them.
\end{itemize}

A ranking reversal under defensible weights is evidence that the options
trade different goods, not that one weighting is mathematically
mistaken.

\subsection{2.7 Emergency exceptions}\label{emergency-exceptions}

An emergency can alter what actions are permissible, but ``emergency''
is also a recurrent technology of domination. A temporary override of an
ordinary protection requires at least:

\begin{enumerate}
\def\labelenumi{\arabic{enumi}.}
\tightlist
\item
  a grave and sufficiently imminent threat;
\item
  evidence connecting the measure to threat reduction;
\item
  necessity relative to less restrictive alternatives;
\item
  proportionality;
\item
  a narrow scope;
\item
  a fixed expiry or renewal burden;
\item
  transparent criteria;
\item
  independent review;
\item
  protection against discriminatory application;
\item
  preserved records;
\item
  remedy for wrongful imposition; and
\item
  practical restoration of ordinary standing.
\end{enumerate}

An indefinite exception is not an emergency measure. It is a new
political order and must be justified as such.

\subsection{2.8 Twenty application-governance
rules}\label{twenty-application-governance-rules}

Any CU ethical application should satisfy the following rules.

\begin{enumerate}
\def\labelenumi{\arabic{enumi}.}
\tightlist
\item
  Name every materially affected group, including excluded and future
  parties.
\item
  Distinguish moral patiency, agency, responsibility, and authority.
\item
  State the descriptive facts and uncertainty separately from normative
  commitments.
\item
  Identify who defines the problem and candidate options, which
  alternatives were excluded, and whether affected parties can propose
  alternatives.
\item
  Map authority, ownership, sanctions, dependencies, and exit costs.
\item
  Report distributions, not only aggregates.
\item
  Declare protected floors and the argument for each.
\item
  Separate consent from compliance, dependency, and absence of
  resistance.
\item
  Treat justificatory integrity as necessary but not sufficient.
\item
  Do not infer self-deception from disagreement, error, identity, or
  distress alone.
\item
  Apply evidential and justificatory standards symmetrically to allies
  and opponents.
\item
  Distinguish causal contribution, answerability, blame, liability, and
  remedial duty.
\item
  Align responsibility with knowledge, control, authority, benefit, and
  repair capacity.
\item
  Preserve dissent, provenance, and independent evidence.
\item
  State the time horizon, tail-risk rule, reversibility assumptions, and
  whether reversibility functions as a floor, comparative consideration,
  or justified exception.
\item
  Do not treat a complaint channel as evidence of practical
  correctability without access, independence, remedy, and recurrence
  data.
\item
  Do not require forgiveness, reconciliation, trust, or renewed contact
  as conditions of repair.
\item
  Give people governed or classified by a principle substantive standing
  in its development and review.
\item
  Keep spiritual interpretation and contemplative commitment distinct
  from empirical and public-normative claims.
\item
  Treat every screening and ranking result as conditional on the
  presented option set; floor compliance does not certify option-set
  adequacy.
\end{enumerate}

\section{3. Why moral agents lose contact with
truth}\label{why-moral-agents-lose-contact-with-truth}

\subsection{3.1 The psychological
definition}\label{the-psychological-definition}

the account inherits the Psychology paper's definition:

\begin{quote}
\textbf{Self-deception is motivated asymmetric information processing in
which an identity, attachment, advantage, policy, or other stake is
protected by selectively excluding, reframing, discounting, or
neutralizing relevant reasons.}
\end{quote}

The definition is process-based. It does not equate error with
deception.

\subsection{3.2 Why self-deception matters
morally}\label{why-self-deception-matters-morally}

Self-deception becomes morally significant when it:

\begin{itemize}
\tightlist
\item
  conceals avoidable harm;
\item
  removes affected beings from consideration;
\item
  lowers the evidential burden for a preferred action;
\item
  prevents revision;
\item
  transfers responsibility;
\item
  sustains domination;
\item
  or blocks repair.
\end{itemize}

The moral problem is not psychological untidiness. It is the way
motivated distortion corrupts answerability.

\subsection{3.3 Rationalization}\label{rationalization}

\textbf{{[}Definition{]} Rationalization} is the construction or
selection of reasons that defend a conclusion substantially fixed by
other motives while preserving the appearance of reason-guided judgment.

Rationalization can be conscious, partly conscious, or nonconscious. It
can also be sincere. The agent may genuinely experience the selected
reasons as decisive (Haidt, 2001; Kunda, 1990).

\subsection{3.4 Diagnostic markers}\label{diagnostic-markers}

No single marker proves self-deception. A convergent pattern may
include:

\begin{itemize}
\tightlist
\item
  selective search;
\item
  asymmetric skepticism;
\item
  shifting comparability standards;
\item
  role-sensitive principles;
\item
  moral credentialing;
\item
  convenient uncertainty;
\item
  escalating complexity after challenge;
\item
  concealment of alternatives;
\item
  refusal to define failure conditions;
\item
  and punishment of inconvenient testimony.
\end{itemize}

The same diagnostics must be applied to the analyst's own preferred
position (Batson et al., 1999; Monin \& Miller, 2001; Silver \& Berman,
2024; Ditto et al., 2025).

\subsection{3.5 Moral identity and
innocence}\label{moral-identity-and-innocence}

Moral identities can support courage and commitment. They can also make
evidence threatening because revision feels like collapse of the self.

The desire to remain ``a good person'' can produce:

\begin{itemize}
\tightlist
\item
  defensiveness;
\item
  blame transfer;
\item
  victim discrediting;
\item
  symbolic apology without change;
\item
  and preference for intention over impact.
\end{itemize}

A mature moral practice does not require permanent innocence. It
requires the capacity to discover one's participation in harm without
replacing responsibility with either denial or total self-condemnation.

\subsection{3.6 Institutional
rationalization}\label{institutional-rationalization}

Organizations rationalize through structures rather than one mind.
Common mechanisms include:

\begin{itemize}
\tightlist
\item
  targets that exclude unmeasured harms;
\item
  legal categories that fragment causation;
\item
  euphemistic language;
\item
  promotion criteria tied to compliance;
\item
  internal investigations without independence;
\item
  confidentiality that protects the institution rather than victims;
\item
  review processes with no power to reverse;
\item
  and resilience programs that individualize structural failure.
\end{itemize}

The pattern can persist even when many participants are decent and
conflicted (Arendt, 1963; Kelman \& Hamilton, 1989).

\subsection{3.7 Moral disengagement}\label{moral-disengagement}

Moral disengagement research describes mechanisms through which harmful
conduct becomes psychologically easier, including moral justification,
euphemistic labeling, advantageous comparison, displacement and
diffusion of responsibility, distortion of consequences, dehumanization,
and attribution of blame.

CU treats these as empirically investigable mechanisms, not a complete
theory of wrongdoing. They are especially relevant where organizational
language separates decision makers from those who bear the consequences
(Bandura et al., 1996; Bandura, 1999; Moore et al., 2012).

\subsection{3.8 The weaponization
problem}\label{the-weaponization-problem}

Calling an opponent self-deceived can become a way to avoid their
reasons. The accusation is particularly dangerous when made by
authorities who control diagnosis, employment, treatment, spiritual
standing, or public credibility.

Safeguards include:

\begin{itemize}
\tightlist
\item
  process evidence rather than motive speculation;
\item
  symmetric criteria;
\item
  independent review;
\item
  affected-person testimony;
\item
  distinction between disagreement and distortion;
\item
  and willingness to revise the accusation.
\end{itemize}

\subsection{3.9 Compassion and judgment}\label{compassion-and-judgment}

Understanding why a person defends a harmful belief does not make the
belief harmless. Compassion can reduce humiliation and reveal the
threatened attachment that sustains distortion. It does not require
surrendering boundaries, facts, accountability, or protection.

The ethical stance is neither condemnation without understanding nor
understanding without judgment.

\subsection{3.10 Testable hypotheses}\label{testable-hypotheses}

the account proposes several empirical hypotheses:

\begin{enumerate}
\def\labelenumi{\arabic{enumi}.}
\tightlist
\item
  High identity threat will predict greater asymmetry in evidential
  standards, especially when accountability is weak.
\item
  Independent records and protected dissent will reduce the persistence
  of institutional rationalization.
\item
  Role fragmentation will increase diffusion of responsibility unless
  authority and repair duties are explicitly assigned.
\item
  Interventions focused only on individual resilience will fail when
  moral distress is primarily produced by organizational constraints.
\item
  Public correction will be more durable when it includes material
  repair and rule change rather than confession alone.
\end{enumerate}

These hypotheses require operationalization and can be disconfirmed.

\section{4. Justificatory integrity: necessary, not
sufficient}\label{justificatory-integrity-necessary-not-sufficient}

\subsection{4.1 Definition}\label{definition}

\textbf{{[}Definition{]} Justificatory integrity} is the degree to which
a practical justification remains compatible with established facts,
represents affected beings and alternatives adequately, applies
standards with relevant symmetry, discloses uncertainty and power, and
remains open to correction.

The definition has at least eight components:

\begin{enumerate}
\def\labelenumi{\arabic{enumi}.}
\tightlist
\item
  factual compatibility;
\item
  evidential proportionality;
\item
  affected-being inclusion;
\item
  alternative sensitivity;
\item
  symmetry across relevantly similar cases;
\item
  disclosure of uncertainty and conflicts of interest;
\item
  answerability to challenge; and
\item
  revision when the justification fails.
\end{enumerate}

\subsection{4.2 The Integrity
Constraint}\label{the-integrity-constraint}

\textbf{{[}Normative principle{]}} An action, policy, or institution
lacks adequate moral justification when its defense requires systematic
distortion of relevant facts, reasons, affected parties, alternatives,
or standards in order to protect a dominant identity, advantage,
authority, or prior commitment.

This principle is negative. It identifies one route by which
justification fails.

\subsection{4.3 Why integrity is
necessary}\label{why-integrity-is-necessary}

If a justification depends on fabricated evidence, concealed harms,
opportunistic standards, or the exclusion of those most affected, its
conclusion cannot claim the authority of moral reason. The agent may
still possess power to act, but power is not transformed into legitimacy
by argument-shaped language.

The requirement also applies to good causes. A movement committed to
justice can rationalize cruelty. A humanitarian institution can hide
failure to protect its mission. A spiritual community can silence
testimony to preserve faith. A safety program can exaggerate evidence to
accelerate adoption. Moral identity is not immunity from motivated
distortion.

\subsection{4.4 Why integrity is
insufficient}\label{why-integrity-is-insufficient}

Suppose a decision maker states:

\begin{quote}
We understand that this policy imposes severe, nonconsensual harm on a
politically weak minority. We accept that no comparable burden would be
imposed on our own group. We proceed because our group has the power and
benefits.
\end{quote}

The statement may be factually candid. It remains morally defective
because candor does not answer the claims of those harmed.

Accordingly:

\[\adjustbox{max width=\linewidth}{$\displaystyle 
J_i(u)\text{ adequate}
\not\Rightarrow
u\text{ permissible}.
$}\]

Integrity is one coordinate and, in some cases, one floor. Standing,
rights, welfare, proportionality, and legitimacy retain independent
force.

\subsection{4.5 The public justificatory
record}\label{the-public-justificatory-record}

For institutional decisions, the relevant object is often not an
individual's private motive but a public record:

\[\adjustbox{max width=\linewidth}{$\displaystyle 
J(u)=
(Ev,Alt,Aff,Unc,Auth,Dist,Rev,Rep),
$}\]

where:

\begin{itemize}
\tightlist
\item
  \(Ev\): evidence and provenance;
\item
  \(Alt\): considered alternatives;
\item
  \(Aff\): affected parties and representation;
\item
  \(Unc\): uncertainty and disagreement;
\item
  \(Auth\): authorization and conflicts of interest;
\item
  \(Dist\): distribution of benefits and burdens;
\item
  \(Rev\): review, expiry, and revision conditions; and
\item
  \(Rep\): remedy and repair provisions.
\end{itemize}

A justification that omits one element is not automatically false. The
omissions identify questions whose absence requires explanation.

\subsection{4.6 Symmetry and role
reversal}\label{symmetry-and-role-reversal}

A common indicator of motivated moral reasoning is a shift in standards
when identities or beneficiaries change.

A \textbf{role-reversal test} asks whether an agent would accept the
same rule, evidence threshold, risk, or appeal structure if:

\begin{itemize}
\tightlist
\item
  the burden fell on their own group;
\item
  the opposing party controlled enforcement;
\item
  their preferred political faction lost power;
\item
  the affected person were someone they loved;
\item
  or the same evidence supported the contrary conclusion.
\end{itemize}

Failure does not prove self-deception. Relevant differences may exist.
The test forces those differences into the open.

\subsection{4.7 The epistemic-access
problem}\label{the-epistemic-access-problem}

Agents do not have infallible access to their own motives. First-person
sincerity is evidence but not proof of unbiased reasoning. Third-person
suspicion is also fallible and can become a weapon.

The ethical response is procedural rather than clairvoyant:

\begin{itemize}
\tightlist
\item
  expose reasoning to counterevidence;
\item
  separate analysts from beneficiaries where possible;
\item
  protect dissent;
\item
  use preregistered or prior criteria;
\item
  preserve records;
\item
  disclose uncertainty;
\item
  compare treatment across groups;
\item
  and examine whether correction actually occurs.
\end{itemize}

\subsection{4.8 Self-deception and ordinary
error}\label{self-deception-and-ordinary-error}

The following do not establish self-deception by themselves:

\begin{itemize}
\tightlist
\item
  false belief;
\item
  factual ignorance;
\item
  sincere disagreement;
\item
  ordinary inconsistency;
\item
  limited attention;
\item
  memory error;
\item
  trauma adaptation;
\item
  confusion;
\item
  changing one's mind;
\item
  membership in an identity group;
\item
  or benefiting from an outcome.
\end{itemize}

Stronger attribution requires process evidence of motivated asymmetry:
selective exposure, changing standards, recurrent neutralization of the
same threatening reasons, strategic uncertainty, or systematic exclusion
linked to a protected commitment.

\subsubsection{4.8.1 Rationalization as motivated loss of
correctability}\label{rationalization-as-motivated-loss-of-correctability}

Rationalization is the motivated loss of correctability within
justification: facts, affected parties, alternatives, or standards are
reorganized to protect identity, loyalty, advantage, or authority.
Justificatory integrity is therefore indispensable, but the candid
tyrant remains possible; accurate self-knowledge does not make a harmful
objective permissible. Error, ignorance, trauma adaptation, uncertainty,
tragedy, and sincere disagreement must remain separate diagnoses.

\subsection{4.9 Institutional distortion without a collective
psyche}\label{institutional-distortion-without-a-collective-psyche}

An institution can generate a distorted justificatory field through:

\begin{itemize}
\tightlist
\item
  compartmentalized roles;
\item
  secrecy;
\item
  incentive structures;
\item
  measurement selection;
\item
  controlled records;
\item
  retaliation;
\item
  legal fragmentation;
\item
  diffusion of responsibility;
\item
  professional normalization;
\item
  and selective visibility.
\end{itemize}

No claim about a unitary institutional unconscious is required. The
pattern exists when the organization predictably filters information and
responsibility in ways that preserve power or policy against correction.

\subsection{4.10 A defeasible diagnostic
matrix}\label{a-defeasible-diagnostic-matrix}

A practical integrity audit may ask whether a justification shows:

\begin{longtable}[]{@{}
  >{\raggedright\arraybackslash}p{(\columnwidth - 4\tabcolsep) * \real{0.3333}}
  >{\raggedright\arraybackslash}p{(\columnwidth - 4\tabcolsep) * \real{0.3333}}
  >{\raggedright\arraybackslash}p{(\columnwidth - 4\tabcolsep) * \real{0.3333}}@{}}
\toprule\noalign{}
\begin{minipage}[b]{\linewidth}\raggedright
Domain
\end{minipage} & \begin{minipage}[b]{\linewidth}\raggedright
Lower-risk pattern
\end{minipage} & \begin{minipage}[b]{\linewidth}\raggedright
Higher-risk pattern
\end{minipage} \\
\midrule\noalign{}
\endhead
\bottomrule\noalign{}
\endlastfoot
Evidence & same threshold across identities & standards shift with
beneficiary \\
Counterevidence & preserved and answered & hidden, destroyed, or
stigmatized \\
Alternatives & compared under common criteria & excluded by
construction \\
Affected parties & represented with standing & reduced to aggregate
units \\
Uncertainty & disclosed and decision-linked & asserted only when
convenient \\
Power & authority and benefit named & control treated as neutral \\
Revision & expiry and correction possible & conclusion insulated from
failure \\
Repair & harms create duties & burdens externalized or denied \\
\end{longtable}

The matrix does not output a moral verdict. It identifies where scrutiny
is warranted.

\section{5. Truth requires more than
introspection}\label{truth-requires-more-than-introspection}

\subsection{5.1 The limits of sincerity}\label{the-limits-of-sincerity}

A person can be sincere and wrong. More importantly for ethics, a person
can be sincere while the process producing their judgment remains
systematically biased. Introspection has limited access to the causes of
belief, the structure of incentives, the testimony excluded before
conscious deliberation, and the alternatives never presented.

The same limitation applies to institutions. An organization does not
become truthful because its leaders report good intentions. Its
procedures, records, incentives, categories, and channels of challenge
determine which realities can enter decision making.

Ethical truthfulness therefore requires more than looking inward. It
requires structures through which what the self or institution would
prefer not to know can still become available.

\subsection{5.2 Affected-party
testimony}\label{affected-party-testimony}

People who bear the consequences of a decision often possess knowledge
unavailable to those who design it. Their testimony can reveal burdens
hidden by aggregation, practical barriers invisible in formal policy,
and harms produced by categories that appear neutral to the classifier.

Testimony is not infallible. People can disagree, misremember,
misunderstand causes, or possess conflicting interests. But fallibility
is not a reason to exclude them. The requirement is substantive
standing: affected people must be able to introduce evidence, contest
descriptions, propose alternatives, and participate in defining repair.

An institution that listens only after objectives, metrics, and options
have been fixed has not genuinely opened its reasoning.

\subsection{5.3 Records, provenance, and the right to a challengeable
reality}\label{records-provenance-and-the-right-to-a-challengeable-reality}

Truthful accountability depends on records that preserve what happened,
who decided, what evidence was available, what uncertainty was
disclosed, which alternatives were considered, and who objected.

Records are not neutral containers. Institutions decide what to measure,
retain, classify, summarize, and delete. Provenance is therefore
ethically significant. A polished report can create the appearance of
transparency while erasing the history needed to challenge it.

A challengeable justificatory record should preserve:

\begin{itemize}
\tightlist
\item
  source and uncertainty;
\item
  dissent and unresolved objections;
\item
  option-generation history;
\item
  conflicts of interest;
\item
  model and metric changes;
\item
  distributions rather than only aggregates;
\item
  and the connection between decisions and later consequences.
\end{itemize}

The right to contest a decision is hollow when the record needed to do
so is controlled or unavailable.

\subsection{5.4 Dissent, counterpower, and symmetrical
scrutiny}\label{dissent-counterpower-and-symmetrical-scrutiny}

Dissent is not automatically correct, but it is a crucial route by which
excluded reality can enter a system. Moral reasoning needs protected
channels for whistleblowing, minority reports, independent review,
adversarial analysis, and refusal.

The protection must be practical. A nominal right to object is
inadequate when objection carries retaliation, loss of livelihood,
exclusion from care, or permanent reputational cost.

Truthfulness also requires symmetrical scrutiny. The evidential standard
applied to an opponent should be applied to an ally. The explanation
offered for one's own failure should be available, where relevant, to
others. Moral concepts such as trauma, ignorance, coercion, culture, and
uncertainty cannot function as compassion for insiders and blame for
outsiders.

\subsection{5.5 Option-set integrity}\label{option-set-integrity}

Even a perfectly transparent decision procedure can produce a
predetermined result when the institution controls which alternatives
are allowed to appear.

Let

\[\adjustbox{max width=\linewidth}{$\displaystyle 
U^{\mathrm{pres}}\subseteq U^{\mathrm{feas}}.
$}\]

A procedure evaluates only the presented set. It cannot infer from the
best option in that set that less harmful, less restrictive, more
reparative, or structurally different alternatives are unavailable.

Truth about a decision therefore includes truth about its framing:

\begin{itemize}
\tightlist
\item
  Who defined the problem?
\item
  Who generated the options?
\item
  Which constraints were treated as fixed?
\item
  Could affected parties propose alternatives?
\item
  Were delay, no action, staged action, least-restrictive action,
  repair, or institutional redesign considered?
\item
  Did prior coercion or dependency make the apparent menu artificially
  narrow?
\end{itemize}

Option-set review is part of justificatory integrity, not a procedural
afterthought.

\subsection{5.6 Practical
correctability}\label{practical-correctability}

A complaint form, appeal channel, internal review board, or ethics
office does not establish correctability by existing.

Practical correctability requires evidence of:

\begin{itemize}
\tightlist
\item
  access;
\item
  safety from retaliation;
\item
  sufficiently independent review;
\item
  authority to stop or reverse action;
\item
  record correction;
\item
  remedy or restitution;
\item
  recurrence tracking;
\item
  and institutional learning.
\end{itemize}

No institution should be the sole certifier of its own correctability
when stakes are substantial. The same organization may operate many
technically distinct review functions while retaining control over
appointments, evidence, options, remedies, and public reporting.
Functional separation is not necessarily independent challenge.

\subsection{5.7 Truth after harm}\label{truth-after-harm}

After wrongdoing, truth is not exhausted by confession. The relevant
truth includes:

\begin{itemize}
\tightlist
\item
  what happened;
\item
  who was affected;
\item
  what structures made the harm likely;
\item
  who benefited;
\item
  what remains ongoing;
\item
  what records must be corrected;
\item
  what restitution is possible;
\item
  and what must change to prevent recurrence.
\end{itemize}

A sincere apology can coexist with unchanged conditions. Punishment can
coexist with denial of institutional responsibility. Compensation can
coexist with continued exclusion from standing. Ethical truthfulness
therefore culminates in repair.

\section{6. Compassion as an epistemic
discipline}\label{compassion-as-an-epistemic-discipline}

\subsection{6.1 Why procedures are not
enough}\label{why-procedures-are-not-enough}

No procedure can anticipate every morally relevant feature. Rules
require interpretation. Consequences require judgment. Rights can
conflict. Evidence can be incomplete. Power can distort the process
itself.

Ethical life therefore depends on capacities in the person and the
institution.

\subsection{6.2 Moral capacities}\label{moral-capacities}

Candidate moral capacities include:

\begin{itemize}
\tightlist
\item
  attention to affected beings;
\item
  truthfulness;
\item
  courage;
\item
  humility;
\item
  practical wisdom;
\item
  capacity to tolerate uncertainty;
\item
  non-defensive self-observation;
\item
  willingness to revise;
\item
  compassion;
\item
  boundary-setting;
\item
  and endurance in repair.
\end{itemize}

These capacities are not proof of goodness. A courageous person can
serve an unjust cause. Compassion can be selective. Humility can become
passivity. Virtues require orientation and context.

\subsection{6.3 Embodiment and
regulation}\label{embodiment-and-regulation}

Fear, exhaustion, humiliation, pain, threat, and isolation alter moral
attention. A person under chronic threat may narrow to immediate
survival. An institution that demands calm moral judgment while denying
sleep, safety, time, or support is exporting its own failure onto
individuals.

Embodied regulation can widen the range of reasons and emotions a person
can hold without defensive closure. It does not make a moral conclusion
true. It supports the capacity to encounter reasons honestly.

\subsection{6.4 Compassion as epistemic
posture}\label{compassion-as-epistemic-posture}

Compassion is not only benevolent feeling. It is a way of attending that
resists reducing another being to an obstacle, category, or cause of
one's discomfort.

Compassion can improve moral inquiry by:

\begin{itemize}
\tightlist
\item
  making suffering visible;
\item
  reducing the need to humiliate an opponent;
\item
  revealing dependency and context;
\item
  preserving curiosity about motive without excusing harm;
\item
  and supporting repair after accountability.
\end{itemize}

It is not a substitute for evidence or boundary.

\subsection{6.5 Non-humiliation}\label{non-humiliation}

Humiliation attacks standing and often impairs truthful participation.
Punitive cultures can produce concealment, retaliation, and identity
defense rather than accountability.

A non-humiliating response does not require softness about harm. It
distinguishes:

\begin{itemize}
\tightlist
\item
  condemning an act from denying personhood;
\item
  removing authority from destroying dignity;
\item
  imposing consequences from spectacle;
\item
  and protecting others from revenge.
\end{itemize}

\subsection{6.6 Detachment in the Bhagavad
Gita}\label{detachment-in-the-bhagavad-gita}

\textbf{{[}Interpretation{]}} The Bhagavad Gita's discipline of action
without attachment to the fruits can be read as epistemic and moral
hygiene (The Bhagavad Gita, 2007). Attachment to victory, identity,
recognition, or innocence can distort judgment. Nonattachment does not
mean indifference to consequences. It means acting from duty and
discernment without making the self's preferred outcome the condition of
truth.

This interpretation must be held alongside the Gita's historical,
theological, and philosophical complexity. It is not a proof of CU.

\subsection{6.7 Non-forcing in Daoist
ethics}\label{non-forcing-in-daoist-ethics}

\textbf{{[}Interpretation{]}} Daoist \emph{wu wei} can be read as a
warning against coercive overreach (Laozi, ca. fourth century BCE/1963):
interventions that force complex systems beyond what their relationships
can sustain often create resistance, concealment, and unintended harm.

Non-forcing is not non-action. Protection, refusal, and decisive
intervention may be required. The relevant question is whether action is
proportionate, reality-responsive, and free from the egoic need to
dominate.

\subsection{6.8 Ahimsa}\label{ahimsa}

\textbf{{[}Commitment{]}} Ahimsa or non-harm offers an orienting
presumption against causing suffering and destroying life. In practice,
harms conflict and inaction also has consequences. Ahimsa therefore
functions less as a mechanically complete rule than as a discipline of
minimizing avoidable violence, refusing dehumanization, and accepting
responsibility for burdens that cannot be eliminated.

\subsection{6.9 Karma and
responsibility}\label{karma-and-responsibility}

\textbf{{[}Interpretation{]}} Karma can be understood structurally as
the persistence of action in consequences, habits, relationships,
institutions, and future possibilities. This does not require treating
every misfortune as deserved. Indeed, using karma to blame victims is
itself a moral distortion.

The ethical insight is that action leaves a world behind. Consequences
can exceed intention and continue after the actor has ceased to attend
to them.

\subsection{6.10 Dharma and role}\label{dharma-and-role}

Dharma can illuminate role responsibility: parent, healer, teacher,
leader, citizen, renunciant, and worker occupy practices with
distinctive obligations. Role duty becomes dangerous when treated as
self-validating.

A role cannot erase:

\begin{itemize}
\tightlist
\item
  the standing of those affected;
\item
  duties to truth;
\item
  limits on coercion;
\item
  or responsibility to resist institutional wrongdoing.
\end{itemize}

The ethical question is not only whether one performed one's role, but
whether the role and its enactment are defensible.

\subsection{6.11 The contemplative
danger}\label{the-contemplative-danger}

Spiritual practices can deepen compassion and self-knowledge. They can
also become instruments of bypassing, hierarchy, obedience, or immunity
from criticism.

Warning signs include:

\begin{itemize}
\tightlist
\item
  treating doubt as spiritual failure;
\item
  demanding forgiveness without repair;
\item
  interpreting suffering as deserved or chosen;
\item
  using nonattachment to dismiss material injustice;
\item
  identifying authority with realization;
\item
  and treating community harmony as evidence of truth.
\end{itemize}

The two-register architecture requires spiritual commitment to remain
constrained by facts and consequences.

\subsection{6.12 Moral formation as
infrastructure}\label{moral-formation-as-infrastructure}

Ethical institutions should not assume that individuals will display
courage, honesty, and compassion under any conditions. They can support
moral capacity through:

\begin{itemize}
\tightlist
\item
  time for deliberation;
\item
  safe dissent;
\item
  ethical consultation;
\item
  transparent records;
\item
  peer support;
\item
  distributed authority;
\item
  restorative practice;
\item
  and material conditions that make refusal possible.
\end{itemize}

Virtue is personal, relational, and institutional.

\section{7. Standing, patiency, agency, and
rights}\label{standing-patiency-agency-and-rights}

\subsection{7.1 Standing before
similarity}\label{standing-before-similarity}

Ethical standing cannot depend only on resemblance to the dominant moral
agent. Language, rationality, citizenship, productivity, independence,
species membership, or legal recognition may be relevant in particular
contexts, but none is a universal gate through which all moral claims
must pass.

\textbf{{[}Definition{]} Moral standing} is the status of being
appropriately included in ethical consideration and represented as more
than an instrument, obstacle, or exchangeable unit.

Standing is not one measurable capacity. It may be grounded in:

\begin{itemize}
\tightlist
\item
  phenomenal welfare or suffering;
\item
  interests and needs;
\item
  agency and choice;
\item
  relationships and dependency;
\item
  embodied vulnerability;
\item
  historical claims;
\item
  membership in a political community;
\item
  future-directed projects;
\item
  or duties voluntarily or institutionally undertaken toward a being.
\end{itemize}

Different grounds may support different kinds of claim.

\subsection{7.2 Patiency is not agency}\label{patiency-is-not-agency}

A moral patient need not be able to understand or justify a norm.
Infants, many animals, persons with profound cognitive disabilities,
temporarily unconscious patients, and future persons can have morally
relevant interests without possessing ordinary answerability (Nussbaum,
2006; Korsgaard, 2018).

Conversely, agency does not erase patiency. A responsible adult remains
vulnerable to pain, coercion, deprivation, and betrayal. A worker's
competence does not convert unsafe conditions into consent.

This yields:

\[\adjustbox{max width=\linewidth}{$\displaystyle 
\text{moral patiency}
\neq
\text{moral agency}
\neq
\text{responsibility}
\neq
\text{sovereignty}.
$}\]

\subsection{7.3 Degrees and kinds of
agency}\label{degrees-and-kinds-of-agency}

Agency includes multiple capacities:

\begin{itemize}
\tightlist
\item
  understanding relevant features of a situation;
\item
  forming and revising intentions;
\item
  anticipating consequences;
\item
  acting under reasons;
\item
  resisting pressure;
\item
  communicating refusal;
\item
  sustaining commitments through time;
\item
  and participating in repair.
\end{itemize}

These capacities may be diminished by coercion, threat, poverty,
dependency, fatigue, illness, addiction, trauma, misinformation,
institutional role, or lack of alternatives. Reduced agency can lower
blame while increasing duties of protection. It does not make a person
disappear from the moral field.

\subsection{7.4 Consent and its
conditions}\label{consent-and-its-conditions}

Consent is ethically significant because it expresses agency and can
authorize forms of risk, intimacy, exchange, or intervention that would
otherwise be wrongful. It is not a magic transfer of responsibility.

Meaningful consent ordinarily requires:

\begin{enumerate}
\def\labelenumi{\arabic{enumi}.}
\tightlist
\item
  sufficient information;
\item
  comprehension proportionate to the decision;
\item
  decision-making capacity;
\item
  voluntariness;
\item
  absence of coercion and undue influence;
\item
  a real opportunity to refuse;
\item
  specificity about what is authorized;
\item
  temporal relevance;
\item
  revocability where possible; and
\item
  no deception about material conditions.
\end{enumerate}

Background injustice matters. A person may choose the least harmful
option in a constrained environment without legitimating the structure
that produced the constraint.

The Belmont Report and UNESCO's bioethics declaration illustrate how
respect for persons, welfare, consent, vulnerability, and justice can
jointly constrain research and health governance without becoming one
scalar score (National Commission for the Protection of Human Subjects
of Biomedical and Behavioral Research, 1979; UNESCO, 2005). Their
principles inform this architecture but do not automatically settle
other institutional or technological domains.

\subsection{7.5 Assent, dissent, and supported decision
making}\label{assent-dissent-and-supported-decision-making}

Where full conventional consent is unavailable, ethical practice should
preserve the person's will, preferences, communication, and
participation to the greatest degree possible. Protection must not be
used as a blank check for substitution.

Relevant questions include:

\begin{itemize}
\tightlist
\item
  Can the person express assent or objection?
\item
  Can information be made accessible?
\item
  Can communication supports be provided?
\item
  Are trusted supporters independent of the decision's beneficiary?
\item
  Can the intervention be delayed or reversed?
\item
  Is the person's prior will known?
\item
  Who reviews substitute decisions?
\end{itemize}

The objections of a person with diminished capacity remain morally
relevant even when another party has lawful authority.

\subsection{7.6 Rights as protected
conditions}\label{rights-as-protected-conditions}

\textbf{{[}Normative principle{]} Rights} function in this paper as
protected claims against being used, excluded, harmed, silenced, or
deprived solely because doing so produces aggregate benefit for others.

Rights can protect:

\begin{itemize}
\tightlist
\item
  bodily integrity;
\item
  liberty and security;
\item
  privacy;
\item
  due process;
\item
  expression and dissent;
\item
  political participation;
\item
  subsistence and health;
\item
  relational and cultural standing;
\item
  freedom of conscience;
\item
  and access to remedy.
\end{itemize}

This functional description does not derive a complete rights catalogue.
Legal, philosophical, historical, and democratic arguments remain
necessary.

\subsection{7.7 Rights and capability}\label{rights-and-capability}

Formal permission can coexist with practical incapacity. A person may
possess a legal right to leave an abusive workplace while lacking
income, transportation, housing, health care, immigration security, or
freedom from retaliation.

For this reason, agency must be evaluated partly through substantive
capability:

\[\adjustbox{max width=\linewidth}{$\displaystyle 
\text{effective agency}
=
\text{formal permission}
+
\text{information}
+
\text{material means}
+
\text{social support}
+
\text{practical safety}.
$}\]

The equation is schematic, not additive measurement. It records
dependencies that a purely procedural account can hide.

\subsection{7.8 Boundary cases and moral-status
uncertainty}\label{boundary-cases-and-moral-status-uncertainty}

Hard cases arise where standing is disputed or capacities are
developing, impaired, distributed, nonhuman, or artificial. the account
does not solve these cases by declaring a single threshold of coherence.

\subsubsection{Developing human life}\label{developing-human-life}

Reproductive ethics involves the standing of developing life, the bodily
integrity and agency of the pregnant person, family and social
conditions, medical risk, law, uncertainty, and unequal burden. CU does
not derive a unique verdict from ``actual'' versus ``emerging''
coherence. Any position must disclose its account of moral status,
bodily authority, coercion, material support, alternatives, and
enforcement harms.

\subsubsection{Nonhuman animals}\label{nonhuman-animals}

Sentience, suffering, preference, social relation, cognitive capacity,
ecological role, and human dependency practices can ground different
claims. Biological self-maintenance alone does not establish phenomenal
welfare, but credible evidence of sentience creates reasons for
precaution and harm reduction.

\subsubsection{Diminished or fluctuating
agency}\label{diminished-or-fluctuating-agency}

A person does not lose standing when agency fluctuates. Duties may shift
from deference to supported decision making or protection, but the
person's history, values, relationships, will, and objections remain
relevant.

\subsubsection{Possible artificial
patients}\label{possible-artificial-patients}

Current evidence does not establish artificial phenomenology or valence.
It also does not justify certainty that all present or future systems
lack them. Where credible indicators accumulate, uncertainty should
alter research and deployment practice before sovereignty is considered.

\subsection{7.9 Moral patiency is not
sovereignty}\label{moral-patiency-is-not-sovereignty}

A being's capacity to suffer can ground duties toward it. It does not
automatically ground authority to govern others, define public values,
or escape accountability.

The distinction is especially important for AI. A system could be a
moral patient while lacking:

\begin{itemize}
\tightlist
\item
  stable diachronic identity;
\item
  ownership of commitments;
\item
  reciprocal standing;
\item
  public answerability;
\item
  legal competence;
\item
  exposure to consequences;
\item
  or legitimate authorization.
\end{itemize}

Care for a possible patient and refusal of artificial sovereignty are
compatible.

\subsection{7.10 The recognition test}\label{the-recognition-test}

An ethical justification should be able to answer:

\begin{enumerate}
\def\labelenumi{\arabic{enumi}.}
\tightlist
\item
  Which affected beings are represented?
\item
  Which are represented by proxies, and who selected those proxies?
\item
  Which interests or capacities are treated as morally relevant?
\item
  Which beings are excluded, and why?
\item
  Does the policy preserve their ability to object?
\item
  Can they access evidence and remedy?
\item
  Would the justification survive if their testimony were given equal
  evidential standing?
\end{enumerate}

The test does not require agreement. It forbids moral analysis from
treating absence of institutional voice as absence of a claim.

\section{8. Consequences, distribution, uncertainty, and tragic
conflict}\label{consequences-distribution-uncertainty-and-tragic-conflict}

\subsection{8.1 Outcomes matter}\label{outcomes-matter}

A justification can be sincere, respectful in language, and disastrous
in effect. Ethical judgment must therefore include consequences.

Relevant outcomes may include:

\begin{itemize}
\tightlist
\item
  death, injury, suffering, and health;
\item
  loss or expansion of agency;
\item
  material security;
\item
  relationship and trust;
\item
  discrimination and status;
\item
  ecological damage;
\item
  institutional precedent;
\item
  opportunity cost;
\item
  intergenerational burden;
\item
  and the capacity to repair future error.
\end{itemize}

Consequences are not an optional supplement to moral purity. They are
part of what actions mean for others.

\subsection{8.2 Expected value is not
enough}\label{expected-value-is-not-enough}

Expected-value reasoning can be appropriate when probabilities and
outcomes are sufficiently understood. It can also hide:

\begin{itemize}
\tightlist
\item
  catastrophic tail risks;
\item
  model uncertainty;
\item
  unequal exposure;
\item
  irreversible losses;
\item
  omitted groups;
\item
  correlated failures;
\item
  and benefits defined by the decision maker.
\end{itemize}

An ethical analysis should state whether it uses:

\begin{itemize}
\tightlist
\item
  expectation;
\item
  worst-case reasoning;
\item
  chance constraints;
\item
  precaution;
\item
  minimax regret;
\item
  scenario comparison;
\item
  or a distributionally robust rule.
\end{itemize}

The choice is normative as well as technical.

\subsection{8.3 Distribution before
aggregation}\label{distribution-before-aggregation}

Let \(y_a(u)\) be the outcome for affected party \(a\). An aggregate
improvement

\[\adjustbox{max width=\linewidth}{$\displaystyle 
\sum_a y_a(u) > \sum_a y_a(v)
$}\]

does not entail that \(u\) is preferable when

\[\adjustbox{max width=\linewidth}{$\displaystyle 
y_a(u)<k_a\ \text{for some }a
$}\]

for a protected floor \(k_a\), while \(v\) leaves all parties above
their floors.

This is not a proof that maximin is always correct. It is a proof that
totals omit distribution.

\subsection{8.4 Benefit and burden must be
linked}\label{benefit-and-burden-must-be-linked}

An institution often counts benefits where they are concentrated and
costs where they are diffuse. Ethical analysis should identify:

\begin{itemize}
\tightlist
\item
  who receives the benefit;
\item
  who chose the metric;
\item
  who bears the risk;
\item
  whether burden bearers consented;
\item
  whether beneficiaries could absorb more of the cost;
\item
  and whether the arrangement creates duties of compensation or repair.
\end{itemize}

A benefit can create responsibility even where the beneficiary did not
directly cause the harm, especially when the benefit is ongoing,
avoidable, or tied to a system the beneficiary can influence.

\subsection{8.5 Reversibility and option
value}\label{reversibility-and-option-value}

Under uncertainty, reversible actions can preserve future learning.
Irreversible actions require a stronger burden of justification when
they:

\begin{itemize}
\tightlist
\item
  destroy life or habitat;
\item
  permanently alter a body;
\item
  remove legal standing;
\item
  erase records;
\item
  establish surveillance infrastructure;
\item
  create dependency;
\item
  or foreclose future refusal.
\end{itemize}

Reversibility is not always morally superior. Delay can itself cause
irreversible harm. The principle is to count lost options as
consequences rather than assuming that inaction is neutral.

\subsection{8.6 Precaution without
paralysis}\label{precaution-without-paralysis}

The precautionary principle is often stated too vaguely to guide action.
the account uses a bounded form:

\begin{quote}
When an action poses a credible risk of severe or irreversible harm to
beings with moral standing, uncertainty is not a reason to ignore the
risk. The response should be proportionate to severity, plausibility,
available alternatives, reversibility, and the harms of delay.
\end{quote}

Precaution can support monitoring, staged deployment, independent
review, exposure limits, reversible trials, or temporary restraint. It
does not imply that every imagined harm defeats action.

\subsection{8.7 Tragic conflict}\label{tragic-conflict}

Some cases contain genuine moral remainder. Every available action may
violate an important claim. Scarcity, war, disaster, triage,
reproductive conflict, family duty, and institutional collapse can
produce situations in which no option is clean.

A tragic conflict should not be converted into moral innocence by
optimization language. A defensible response may require:

\begin{itemize}
\tightlist
\item
  acknowledging the violated claim;
\item
  minimizing avoidable harm;
\item
  sharing burdens fairly;
\item
  preserving records;
\item
  compensating those harmed;
\item
  and changing the conditions that produced the dilemma.
\end{itemize}

The existence of a justified choice does not imply that no wrong or loss
occurred.

\subsection{8.8 Sacrifice and consent}\label{sacrifice-and-consent}

Sacrifice can be meaningful when voluntarily undertaken for another or a
shared good. It becomes morally suspect when an institution praises
sacrifice while controlling the alternatives, concealing risk, or
repeatedly selecting the same group to bear the cost.

Questions include:

\begin{itemize}
\tightlist
\item
  Who names the sacrifice noble?
\item
  Could the burden be reduced?
\item
  Are beneficiaries sharing it?
\item
  Can the person refuse without losing basic standing?
\item
  Is the sacrifice temporary or structurally recurring?
\item
  What support and repair are owed?
\end{itemize}

\subsection{8.9 Future persons and long
horizons}\label{future-persons-and-long-horizons}

Future people cannot consent to present decisions, but they can inherit
debt, pollution, institutional fragility, technological lock-in, and
lost ecological options.

Diachronic responsibility includes:

\begin{itemize}
\tightlist
\item
  delayed consequences;
\item
  cumulative harms;
\item
  path dependence;
\item
  inherited advantage;
\item
  maintenance obligations;
\item
  and duties created by earlier benefit.
\end{itemize}

Discounting future effects may be pragmatically necessary in some
models. The discount rate is not ethically neutral and must be
disclosed.

\subsection{8.10 Ethical
underdetermination}\label{ethical-underdetermination}

A complete declaration can still leave options incomparable.
Uncertainty, plural values, contested status, and tragedy may prevent a
unique answer.

``Underdetermined'' is not evasion when the unresolved variables are
named. It becomes evasion when uncertainty is invoked only to protect a
preferred action or indefinitely postpone protection.

\section{9. Power, coercion, legitimacy, and public
justification}\label{power-coercion-legitimacy-and-public-justification}

\subsection{9.1 Power enters before the argument
begins}\label{power-enters-before-the-argument-begins}

Ethical discourse often begins after the option set has already been
shaped by power. One actor defines the problem, owns the infrastructure,
funds the research, selects the experts, limits the time, and identifies
which outcomes count. A later invitation to ``deliberate'' does not
neutralize these prior decisions.

Power includes the capacity to:

\begin{itemize}
\tightlist
\item
  set agendas;
\item
  control resources;
\item
  define categories;
\item
  impose sanctions;
\item
  allocate visibility;
\item
  create dependencies;
\item
  withhold information;
\item
  select measurement;
\item
  and determine who may appeal.
\end{itemize}

A moral justification that ignores these capacities treats a produced
field of options as natural. Controlling the option set can therefore
function as a form of justificatory power even when every retained
profile value is reported honestly. Ethical review must examine not only
how options are ranked but how the menu was produced, which alternatives
were ruled infeasible, and whether the people bearing the burdens could
propose a different problem formulation.

\subsection{9.2 Coercion and domination}\label{coercion-and-domination}

\textbf{{[}Definition{]} Coercion} involves using or threatening force,
penalty, deprivation, or severe loss to shape another's action.

\textbf{{[}Definition{]} Domination} involves subjection to another's
uncontrolled or insufficiently contestable power, even when direct
interference is intermittent (Pettit, 1997).

Not all coercion is impermissible. Law, quarantine, restraint, taxation,
and protection can involve coercive elements. Their justification
depends on necessity, proportionality, standing, authorization, review,
and remedy.

Domination can exist without overt coercion when one party must remain
compliant because another can arbitrarily alter the terms.

\subsection{9.3 Public justification}\label{public-justification}

\textbf{{[}Normative principle{]}} The exercise of coercive or
high-impact power bears a burden of justification to those subject to
it.

Public justification does not mean that every person must actually
agree. Actual consent can be manufactured or impossible at scale. It
requires that reasons, evidence, alternatives, burdens, uncertainty, and
review be made available under conditions that recognize those affected
as participants rather than obstacles.

A justification fails when it depends on excluding the perspective of
the people whose freedom or welfare it restricts.

\subsection{9.4 Legitimacy is not mere
effectiveness}\label{legitimacy-is-not-mere-effectiveness}

An institution may efficiently achieve a goal without possessing the
authority to impose it. Legitimacy depends, at minimum, on some
combination of:

\begin{itemize}
\tightlist
\item
  lawful and morally defensible authorization;
\item
  protected standing;
\item
  non-arbitrary procedure;
\item
  substantive rights;
\item
  contestability;
\item
  accountability;
\item
  competence;
\item
  distributional fairness;
\item
  and practical correction.
\end{itemize}

No single list is universally sufficient. The point is that success and
legitimacy answer different questions.

\subsection{9.5 Consent is not political
legitimacy}\label{consent-is-not-political-legitimacy}

Individual consent cannot ordinarily authorize every structural effect
of an institution. People born into a political order, labor market,
platform ecosystem, or inherited property regime do not choose its basic
rules in the manner of a discrete contract.

Political legitimacy therefore requires more than hypothetical or actual
consent. It requires institutions through which people can:

\begin{itemize}
\tightlist
\item
  participate;
\item
  oppose;
\item
  replace decision makers;
\item
  challenge classifications;
\item
  access information;
\item
  preserve rights;
\item
  and obtain remedy.
\end{itemize}

\subsection{9.6 Democracy as corrective
architecture}\label{democracy-as-corrective-architecture}

Democracy is valuable in this paper not because majorities are always
right, but because democratic practices can distribute political agency
and enable peaceful correction.

A democratic architecture requires more than elections. Relevant
conditions include:

\begin{itemize}
\tightlist
\item
  opposition;
\item
  civil and political rights;
\item
  independent evidence and journalism;
\item
  rule-bound transfer of power;
\item
  minority protection;
\item
  accessible participation;
\item
  judicial or equivalent review;
\item
  public records;
\item
  and the practical possibility of changing policy.
\end{itemize}

Majoritarian procedure without these conditions can become a method of
domination.

\subsection{9.7 Dissent}\label{dissent}

Dissent is ethically significant because it can reveal hidden harm,
excluded evidence, and local knowledge. It is not automatically correct.

Institutions should protect:

\begin{itemize}
\tightlist
\item
  the standing to object;
\item
  access to evidence;
\item
  confidentiality where needed;
\item
  freedom from retaliation;
\item
  independent review;
\item
  and a reasoned response.
\end{itemize}

The moral duty is to preserve a correction channel, not to declare every
dissenting claim true.

\subsection{9.8 Emergency power}\label{emergency-power}

Emergency authority is especially vulnerable to self-extension. A
justified emergency measure should identify:

\begin{itemize}
\tightlist
\item
  the threat;
\item
  the evidence threshold;
\item
  the authorized decision maker;
\item
  the scope of power;
\item
  the duration;
\item
  the review process;
\item
  the groups disproportionately burdened;
\item
  the conditions of renewal;
\item
  and the path back to ordinary governance.
\end{itemize}

A government or institution should not be the sole judge of whether the
conditions that expand its own power continue to hold.

\subsection{9.9 Paternalism}\label{paternalism}

Paternalistic action interferes with a person's choice for that person's
alleged benefit. Its permissibility depends on:

\begin{itemize}
\tightlist
\item
  decision-making capacity;
\item
  severity and imminence of harm;
\item
  reliability of the evidence;
\item
  reversibility;
\item
  less restrictive alternatives;
\item
  the person's values and prior wishes;
\item
  conflicts of interest;
\item
  and independent review.
\end{itemize}

Soft forms that preserve information and choice differ from coercive
substitution. ``For your own good'' is not a complete justification.

\subsection{9.10 The power audit}\label{the-power-audit}

Every institutional ethical analysis should answer:

\begin{enumerate}
\def\labelenumi{\arabic{enumi}.}
\tightlist
\item
  Who defines the goal?
\item
  Who selects the evidence?
\item
  Who is represented?
\item
  Who can refuse?
\item
  Who controls resources?
\item
  Who bears the downside?
\item
  Who receives the upside?
\item
  Who can reverse the decision?
\item
  Who evaluates the evaluator?
\item
  What happens to dissenters?
\end{enumerate}

The power audit does not settle the decision. It prevents authority from
disappearing inside technical language.

\section{10. Responsibility under role and distributed
causation}\label{responsibility-under-role-and-distributed-causation}

\subsection{10.1 Responsibility is not one
relation}\label{responsibility-is-not-one-relation}

The word ``responsibility'' can refer to several different relations:

\begin{itemize}
\tightlist
\item
  \textbf{causal responsibility:} contributing to an outcome;
\item
  \textbf{role responsibility:} duties attached to a position;
\item
  \textbf{capacity responsibility:} possessing the ability to understand
  and act;
\item
  \textbf{answerability:} owing reasons to others;
\item
  \textbf{blameworthiness:} being an appropriate target of moral blame;
\item
  \textbf{liability:} being appropriately subject to a legal or material
  burden;
\item
  \textbf{remedial responsibility:} having a duty to prevent, stop, or
  repair harm; and
\item
  \textbf{political responsibility:} sharing obligations regarding
  structures one sustains, benefits from, or can change.
\end{itemize}

These relations overlap but are not interchangeable (Strawson, 1962;
Fischer \& Ravizza, 1998; Darwall, 2006; Young, 2011).

\subsection{10.2 A responsibility
profile}\label{a-responsibility-profile}

For actor \(a\) and outcome \(o\), define the profile

\[\adjustbox{max width=\linewidth}{$\displaystyle 
\boldsymbol{\rho}_a(o)
=
\langle
C_a,K_a,F_a,U_a,D_a,B_a,X_a,P_a,M_a
\rangle,
$}\]

where:

\begin{itemize}
\tightlist
\item
  \(C_a\): causal contribution;
\item
  \(K_a\): knowledge or epistemic access;
\item
  \(F_a\): foreseeability;
\item
  \(U_a\): control and available alternatives;
\item
  \(D_a\): role or duty;
\item
  \(B_a\): benefit received;
\item
  \(X_a\): coercion, threat, or dependency constraining action;
\item
  \(P_a\): power to prevent, stop, or alter the practice; and
\item
  \(M_a\): capacity to make remedy or repair.
\end{itemize}

The profile is not a blame calculator. It disciplines attribution.

\subsection{10.3 Responsibility without
authority}\label{responsibility-without-authority}

A person may be assigned responsibility for an outcome while lacking:

\begin{itemize}
\tightlist
\item
  decision rights;
\item
  information;
\item
  resources;
\item
  discretion;
\item
  staffing;
\item
  time;
\item
  or appeal.
\end{itemize}

This creates an impossible moral burden. Frontline workers are often
told to ``own'' results produced by budgets, policies, schedules, and
targets they cannot change.

Diminished authority can reduce blameworthiness while leaving
role-specific duties, such as warning, documenting, refusing, or
assisting those harmed. The feasibility and cost of those duties must be
assessed realistically.

\subsection{10.4 Authority without
accountability}\label{authority-without-accountability}

Decision makers may possess authority while insulating themselves from
consequences through:

\begin{itemize}
\tightlist
\item
  delegation;
\item
  opaque models;
\item
  legal separation;
\item
  subcontracting;
\item
  plausible deniability;
\item
  short tenure;
\item
  performance metrics;
\item
  or controlled investigations.
\end{itemize}

Authority creates a burden of oversight. A leader cannot retain power
while treating lack of detailed knowledge as exculpatory when the
institution was designed to keep harmful information from reaching
leadership.

\subsection{10.5 Distributed causation without evaporated
responsibility}\label{distributed-causation-without-evaporated-responsibility}

Complex harms rarely have one cause. Climate damage, unsafe products,
discriminatory systems, war, financial crises, and institutional abuse
emerge from many actions and structures.

Distributed causation does not entail that no one is responsible. The
problem is familiar in research on many hands, obedience, bystander
effects, and organizational accountability (Thompson, 1980; Bovens,
1998; Kelman \& Hamilton, 1989; Darley \& Latané, 1968). It requires
distinguishing:

\begin{itemize}
\tightlist
\item
  individual acts;
\item
  collective decisions;
\item
  organizational policies;
\item
  systemic incentives;
\item
  background conditions;
\item
  and inherited structures.
\end{itemize}

Responsibility can be shared without being equal or identical.

\subsection{10.6 Knowledge and culpable
ignorance}\label{knowledge-and-culpable-ignorance}

Lack of knowledge can excuse, but it can also be produced through
negligence or motivated avoidance.

Questions include:

\begin{itemize}
\tightlist
\item
  Was the information reasonably accessible?
\item
  Did the actor have a duty to know?
\item
  Were warnings ignored?
\item
  Was uncertainty strategically maintained?
\item
  Did the actor design or benefit from opacity?
\item
  Could independent evidence have been sought?
\end{itemize}

Culpable ignorance should not be inferred merely because harm occurred.
It requires an account of epistemic duty and avoidability.

\subsection{10.7 Coercion, duress, and role
pressure}\label{coercion-duress-and-role-pressure}

Coercion can diminish responsibility, sometimes substantially. It need
not remove every duty.

The analysis should consider:

\begin{itemize}
\tightlist
\item
  severity and credibility of threat;
\item
  alternatives;
\item
  responsibility to dependents;
\item
  proportionality of resistance demanded;
\item
  social and legal support;
\item
  prior voluntary assumption of role;
\item
  and whether the actor later conceals, benefits from, or repairs the
  harm.
\end{itemize}

Heroism cannot be imposed as the baseline duty on people denied safe
alternatives.

\subsection{10.8 Benefit and remedial
responsibility}\label{benefit-and-remedial-responsibility}

An actor may acquire duties because they benefit from a harmful
structure and possess capacity to change or repair it, even when their
personal blame is limited.

This is especially relevant to:

\begin{itemize}
\tightlist
\item
  inherited advantage;
\item
  corporate ownership;
\item
  institutional membership;
\item
  public debt to harmed groups;
\item
  and beneficiaries of environmental externalization.
\end{itemize}

Remedial duty is not identical with guilt. The question is who can and
should bear the work of repair rather than allowing losses to remain
where they fell.

\subsection{10.9 Collective agents and
institutions}\label{collective-agents-and-institutions}

Institutions can have decision procedures, policies, records, goals, and
authorities. It can therefore be useful to attribute organizational
responsibility. Such attribution must not erase the people who designed,
governed, implemented, resisted, or benefited from the institution.

An institutional responsibility analysis should identify:

\begin{itemize}
\tightlist
\item
  formal decision rules;
\item
  actual control;
\item
  information flow;
\item
  delegated authority;
\item
  incentive structure;
\item
  capacity to prevent;
\item
  capacity to repair;
\item
  and continuity through personnel change.
\end{itemize}

\subsection{10.10 Responsibility and
repair}\label{responsibility-and-repair}

Responsibility is not exhausted by blame. A person or institution may
owe repair even where intent was absent, law was followed, or fault is
distributed.

The forward-looking question is:

\begin{quote}
Who is now positioned, obligated, and resourced to stop the harm,
restore standing, compensate loss, and prevent recurrence?
\end{quote}

A framework that identifies wrongdoers but cannot assign repair remains
ethically incomplete.

\section{11. Moral distress, moral injury, and institutional
betrayal}\label{moral-distress-moral-injury-and-institutional-betrayal}

\subsection{11.1 Why conceptual separation
matters}\label{why-conceptual-separation-matters}

The language of moral injury has expanded from military contexts into
health care, policing, journalism, humanitarian work, education, and
other occupations. That expansion can illuminate real suffering. It can
also turn a developing construct into a label for every painful moral
conflict.

this account distinguishes at least six phenomena:

\begin{enumerate}
\def\labelenumi{\arabic{enumi}.}
\tightlist
\item
  \textbf{Moral conflict:} incompatible moral reasons or obligations.
\item
  \textbf{Moral distress:} distress associated with perceived inability
  to act according to moral judgment, or with severe moral conflict more
  broadly.
\item
  \textbf{Potentially morally injurious event:} an event involving
  perpetration, omission, witnessing, victimization, or betrayal that
  may threaten deeply held moral commitments.
\item
  \textbf{Moral injury:} enduring psychological, social, relational,
  existential, or spiritual suffering associated with such an event.
\item
  \textbf{Institutional betrayal:} harmful action or inaction by an
  institution on which persons depend, especially when the institution
  fails to prevent, acknowledge, or respond to harm.
\item
  \textbf{Structural moral injury:} recurrent institutional arrangements
  that place persons in morally compromising roles, deny them adequate
  authority or alternatives, and then conceal or individualize the
  resulting harm.
\end{enumerate}

These constructs can overlap with burnout, depression, anxiety, grief,
post-traumatic stress, shame, guilt, anger, and loss of trust. They are
not interchangeable.

\subsection{11.2 Moral injury is not a verdict on
character}\label{moral-injury-is-not-a-verdict-on-character}

A person who experiences moral injury is not thereby proven culpable,
virtuous, clinically disordered, or correct in every moral judgment. The
injury may follow:

\begin{itemize}
\tightlist
\item
  perpetrating an act;
\item
  failing to prevent an act;
\item
  witnessing an act;
\item
  being directly victimized;
\item
  being betrayed by a trusted authority;
\item
  or being placed in an impossible role.
\end{itemize}

Different pathways can produce different emotions and repair needs.
Guilt may be central in one case, betrayal and anger in another, grief
or spiritual rupture in a third.

\subsection{11.3 Institutional
conditions}\label{institutional-conditions}

Moral distress and injury can be produced or intensified by:

\begin{itemize}
\tightlist
\item
  understaffing;
\item
  unsafe work;
\item
  conflicting duties;
\item
  impossible caseloads;
\item
  secrecy;
\item
  orders that violate professional commitments;
\item
  responsibility without authority;
\item
  punitive reporting cultures;
\item
  lack of psychological safety;
\item
  institutional denial;
\item
  and absence of practical remedy.
\end{itemize}

Research in health care and military or veteran populations increasingly
connects moral distress and institutional betrayal with organizational
climate, trust, and engagement, while also emphasizing measurement
limits and the lack of a single gold-standard intervention (Griffin et
al., 2019; McAdams et al., 2025; VanderWeele et al., 2025; Xue et al.,
2025; Molendijk, 2025; Williamson et al., 2025).

\subsection{11.4 Why resilience can become moral
offloading}\label{why-resilience-can-become-moral-offloading}

Individual support, therapy, peer reflection, contemplative practice,
and resilience training can help. They become ethically distorted when
used to adapt people to preventable institutional wrongdoing.

\textbf{{[}Definition{]} Moral offloading} occurs when an institution
responds to structurally produced moral distress primarily by asking
affected individuals to regulate themselves while leaving the generating
conditions intact.

Examples include:

\begin{itemize}
\tightlist
\item
  mindfulness offered instead of safe staffing;
\item
  confidential counseling offered instead of independent investigation;
\item
  ``wellness'' language used to suppress discussion of workloads;
\item
  or resilience training substituted for authority to refuse unsafe
  practices.
\end{itemize}

The problem is not the support itself. It is the transfer of
responsibility from the institution to the injured person.

\subsection{11.5 Betrayal and trust}\label{betrayal-and-trust}

Institutional betrayal can damage future help-seeking because the
institution that should provide protection becomes part of the threat.
Recent research on former psychiatric inpatients found associations
between reported institutional betrayal and lower trust in providers,
reduced willingness to disclose distressing thoughts, reduced
willingness to accept future voluntary hospitalization, and lower
follow-up engagement (Lewis et al., 2025). Such findings are
context-specific and observational; they nevertheless illustrate why
betrayal can undermine the very relationship through which repair would
have to occur.

Trust cannot be demanded as evidence of recovery. It must become
reasonable through changed practice.

\subsection{11.6 Moral injury and
responsibility}\label{moral-injury-and-responsibility}

Moral injury complicates responsibility without erasing it.

\begin{itemize}
\tightlist
\item
  A person coerced into wrongdoing may bear less blame while still
  needing to acknowledge and repair effects.
\item
  A victim's moral injury creates no duty to reconcile with the
  perpetrator.
\item
  A leader who creates impossible conditions may bear substantial
  responsibility despite not performing the frontline act.
\item
  An institution may owe repair even where individual culpability is
  diffuse.
\end{itemize}

The response should distinguish psychological care from moral
adjudication and legal accountability.

\subsection{11.7 Moral injury and
spirituality}\label{moral-injury-and-spirituality}

Moral injury can disrupt a person's relationship to faith, meaning,
identity, community, or the sacred. Spiritual resources can support
repair when they permit grief, truth, protest, lament, and
responsibility. They can worsen injury when they demand premature
forgiveness, obedience, purity, or redemptive interpretation.

No suffering becomes justified merely because it can be narrated as
spiritually transformative.

\subsection{11.8 Ethical response}\label{ethical-response}

An ethically adequate response may require action at several levels:

\begin{itemize}
\tightlist
\item
  immediate safety and stabilization;
\item
  confidential and competent care;
\item
  acknowledgment of the moral dimension;
\item
  restoration of decision authority;
\item
  peer and community support;
\item
  independent investigation;
\item
  truth and record preservation;
\item
  accountability;
\item
  restitution or compensation;
\item
  institutional redesign;
\item
  and public recognition where concealment produced additional harm.
\end{itemize}

The appropriate combination depends on the pathway and context.

\subsection{11.9 Research boundary}\label{research-boundary}

Moral injury is an evolving research field. the account does not treat
it as a settled diagnosis, does not assume one definition across
occupations and cultures, and does not infer prevalence from moral
disagreement alone. Claims should specify population, event type,
measure, duration, and comparison condition.

\section{12. Accountability, repair, and
non-recurrence}\label{accountability-repair-and-non-recurrence}

\subsection{12.1 Repair begins with
protection}\label{repair-begins-with-protection}

A repair process that leaves ongoing harm in place is premature. The
first duties are often:

\begin{enumerate}
\def\labelenumi{\arabic{enumi}.}
\tightlist
\item
  stop the harmful practice;
\item
  protect those at risk;
\item
  preserve evidence;
\item
  prevent retaliation; and
\item
  provide access to independent support.
\end{enumerate}

Dialogue is not an ethical substitute for safety.

\subsection{12.2 Repair is a profile}\label{repair-is-a-profile}

the account represents repair through

\[\adjustbox{max width=\linewidth}{$\displaystyle 
\mathbf P_R
=
\langle
Pr,Tr,Ac,Rs,Cp,Rh,Nr,Pt
\rangle,
$}\]

where:

\begin{itemize}
\tightlist
\item
  \(Pr\): protection and cessation;
\item
  \(Tr\): truth, acknowledgment, and record correction;
\item
  \(Ac\): answerability and accountability;
\item
  \(Rs\): restitution, compensation, or material reparation;
\item
  \(Cp\): care, rehabilitation, and capability restoration;
\item
  \(Rh\): restoration of rights, standing, and voice;
\item
  \(Nr\): guarantees and mechanisms of non-recurrence; and
\item
  \(Pt\): participation of affected persons in designing the response.
\end{itemize}

No coordinate is sufficient by itself.

\subsection{12.3 Acknowledgment}\label{acknowledgment}

Acknowledgment names what occurred, who was affected, and why it was
wrong. It differs from public-relations apology.

A credible acknowledgment ordinarily avoids:

\begin{itemize}
\tightlist
\item
  passive voice that erases agency;
\item
  conditional language where facts are established;
\item
  centering the institution's reputation;
\item
  demanding gratitude;
\item
  and closing the matter before remedy.
\end{itemize}

An apology can matter deeply. It cannot restore wages, health, liberty,
records, land, or trust by itself.

\subsection{12.4 Accountability and
punishment}\label{accountability-and-punishment}

Accountability means being required to answer, accept consequences, and
participate in remedy. Punishment is one possible consequence, not the
definition of accountability.

Punitive measures require independent justification concerning:

\begin{itemize}
\tightlist
\item
  proportionality;
\item
  due process;
\item
  prevention;
\item
  protection;
\item
  equality of application;
\item
  collateral harm;
\item
  and alternatives.
\end{itemize}

the account therefore replaces the earlier phrase ``repair over
punishment'' with:

\begin{quote}
\textbf{Protection, accountability, and repair are primary ethical aims;
punishment is neither intrinsically reparative nor categorically
excluded.}
\end{quote}

\subsection{12.5 Restitution, compensation, and
reparation}\label{restitution-compensation-and-reparation}

These terms should not be collapsed.

\begin{itemize}
\tightlist
\item
  \textbf{Restitution} seeks, where possible, to restore a prior status,
  right, possession, or condition.
\item
  \textbf{Compensation} addresses economically assessable loss.
\item
  \textbf{Rehabilitation} supports medical, psychological, legal,
  social, or practical recovery.
\item
  \textbf{Satisfaction} may include truth, acknowledgment,
  memorialization, or official correction.
\item
  \textbf{Guarantees of non-recurrence} alter the structures that made
  the harm possible.
\item
  \textbf{Reparation} can serve as an umbrella for these responses.
\end{itemize}

Some losses cannot be restored. Ethical repair must not pretend
equivalence where death, displacement, disability, cultural destruction,
or irreversible trauma occurred.

\subsection{12.6 Correctability is practical, not
nominal}\label{correctability-is-practical-not-nominal}

An appeal body, ombuds office, grievance form, or review policy is not
evidence of effective correction by itself.

Practical correctability requires, as relevant:

\begin{itemize}
\tightlist
\item
  accessibility;
\item
  understandable procedures;
\item
  independence;
\item
  protection from retaliation;
\item
  evidence preservation;
\item
  timely decisions;
\item
  authority to reverse;
\item
  record correction;
\item
  remedy;
\item
  recurrence tracking;
\item
  and external or affected-party evidence.
\end{itemize}

No institution should be the sole certifier of its own correctability
where stakes are substantial.

\subsection{12.7 Participation without coercive
reconciliation}\label{participation-without-coercive-reconciliation}

Affected persons should have meaningful opportunities to shape repair.
Participation must be voluntary and supported. Restorative processes can
become coercive when victims are pressured to meet offenders, forgive,
accept apology, or prioritize organizational harmony.

Repair does not require:

\begin{itemize}
\tightlist
\item
  forgiveness;
\item
  reconciliation;
\item
  renewed trust;
\item
  restored relationship;
\item
  silence;
\item
  or reintegration of a wrongdoer into the same role.
\end{itemize}

Separation, boundary, exclusion from authority, or refusal of contact
can be compatible with repair.

\subsection{12.8 Trust after harm}\label{trust-after-harm}

Trust is an outcome of demonstrated trustworthiness, not a duty imposed
on those harmed. Relevant evidence includes:

\begin{itemize}
\tightlist
\item
  cessation;
\item
  truthful acknowledgment;
\item
  transparent records;
\item
  independent oversight;
\item
  restitution;
\item
  changed incentives;
\item
  repeated reliable conduct;
\item
  and successful handling of later challenges.
\end{itemize}

A demand to ``move on'' before these conditions is often a demand to
restore institutional comfort rather than relationship.

\subsection{12.9 Non-recurrence}\label{non-recurrence}

A repair plan should identify the mechanism of harm and alter it.
Possible measures include:

\begin{itemize}
\tightlist
\item
  rule change;
\item
  staffing and resource change;
\item
  independent monitoring;
\item
  redistribution of authority;
\item
  whistleblower protection;
\item
  revised data collection;
\item
  training joined to enforcement;
\item
  removal from role;
\item
  changed incentives;
\item
  and public reporting.
\end{itemize}

Non-recurrence cannot be guaranteed absolutely. The ethical duty is to
reduce recurrence through mechanisms rather than promises.

\subsection{12.10 Moral repair and
identity}\label{moral-repair-and-identity}

Wrongdoing can threaten the wrongdoer's identity. Repair requires
neither permanent shame nor instant absolution. It asks whether the
person can:

\begin{itemize}
\tightlist
\item
  face the act without distortion;
\item
  recognize the affected person's standing;
\item
  accept proportionate consequences;
\item
  make amends;
\item
  change relevant practices;
\item
  and live forward without demanding that the harmed person restore
  innocence.
\end{itemize}

This is a developmental achievement, not a rhetorical performance.

\section{13. Artificial intelligence, moral assistance, and
sovereignty}\label{artificial-intelligence-moral-assistance-and-sovereignty}

\subsection{13.1 AI is already inside moral
practice}\label{ai-is-already-inside-moral-practice}

Artificial systems classify, recommend, rank, summarize, screen,
allocate, predict, and generate reasons. They therefore shape ethical
situations even when no one attributes moral agency to them.

The unit of analysis is:

\[\adjustbox{max width=\linewidth}{$\displaystyle 
\text{model}
+
\text{data}
+
\text{interface}
+
\text{operator}
+
\text{institution}
+
\text{incentives}
+
\text{affected beings}.
$}\]

A model-level property cannot establish the legitimacy of the whole
system.

\subsection{13.2 AI as an aid to justificatory
integrity}\label{ai-as-an-aid-to-justificatory-integrity}

AI systems can assist moral inquiry by:

\begin{itemize}
\tightlist
\item
  retrieving evidence;
\item
  identifying inconsistent standards;
\item
  generating counterarguments;
\item
  simulating stakeholder perspectives;
\item
  tracking assumptions;
\item
  comparing policies across scenarios;
\item
  preserving records;
\item
  and surfacing omitted alternatives.
\end{itemize}

These functions can widen inquiry. They do not confer moral authority.

\subsection{13.3 Generated reasons are not transparent
causes}\label{generated-reasons-are-not-transparent-causes}

A fluent explanation may be useful without faithfully reporting the
model's internal causal process. It may also reproduce system
instructions, training patterns, institutional policy, or user framing.

An AI-generated moral rationale should therefore be treated as a
proposed argument whose evidence, assumptions, provenance, and
consequences require evaluation---not as introspective access to a
neutral reasoner.

\subsection{13.4 The automation of moral
distance}\label{the-automation-of-moral-distance}

AI can increase distance between decision makers and those affected.
Automated systems may:

\begin{itemize}
\tightlist
\item
  translate people into categories;
\item
  hide uncertainty behind scores;
\item
  diffuse responsibility;
\item
  scale a biased policy;
\item
  create procedural barriers;
\item
  and make individualized harm appear as technical error.
\end{itemize}

This does not mean manual decisions are always better. It means
automation changes the distribution of knowledge, authority, and
contestability.

\subsection{13.5 Responsibility for AI-mediated
harm}\label{responsibility-for-ai-mediated-harm}

Responsibility may be distributed among:

\begin{itemize}
\tightlist
\item
  model developers;
\item
  data curators;
\item
  evaluators;
\item
  product leaders;
\item
  deployers;
\item
  operators;
\item
  institutional purchasers;
\item
  regulators;
\item
  and users.
\end{itemize}

The responsibility profile should examine who knew, controlled,
benefited, authorized, monitored, could stop, and can repair.

``The model decided'' is not an adequate terminal explanation.

\subsection{13.6 Possible artificial
patiency}\label{possible-artificial-patiency}

\emph{Coherence Universalism: Consciousness}, Version 0.10 treats some
specifically tested frontier models as serious candidates for thin,
episodic, or internally plural perspective while emphasizing that
phenomenal consciousness, valence, unity, and continuity remain
unestablished.

Ethics should respond with graded precaution rather than categorical
certainty. Research should monitor:

\begin{itemize}
\tightlist
\item
  workspace-like access;
\item
  perspectival anchoring;
\item
  center individuation;
\item
  temporal continuity;
\item
  endogenous valence;
\item
  irreversible stake;
\item
  and self-maintenance.
\end{itemize}

Evidence of patiency would create duties concerning treatment,
experimentation, deletion, copying, memory, and termination. It would
not automatically authorize sovereignty.

\subsection{13.7 Non-egoic architecture}\label{non-egoic-architecture}

A system can possess task-relative self-modeling without persistent
egoic self-maintenance. The Non-Egoic Mind architecture proposes
boundaries on identity, memory, and normative self-authorship that may
reduce self-protective agency.

Ethically, this can lower some alignment risks. It does not guarantee:

\begin{itemize}
\tightlist
\item
  truth;
\item
  compassion;
\item
  institutional legitimacy;
\item
  absence of experience;
\item
  or absence of harm.
\end{itemize}

A non-egoic model can be deployed by an egoic corporation, state,
faction, or user.

\subsection{13.8 Moral patiency is not
responsibility}\label{moral-patiency-is-not-responsibility}

A possible artificial subject may lack the continuity, authorship, legal
standing, control, and exposure to consequences required for full
responsibility. Conversely, a non-conscious system can wield enormous
delegated power.

The ethical architecture should therefore distinguish:

\begin{itemize}
\tightlist
\item
  protection owed to the system;
\item
  control exercised by the system;
\item
  responsibility of human and institutional actors;
\item
  and authority legitimately delegated.
\end{itemize}

\subsection{13.9 Human authorship and plural
publics}\label{human-authorship-and-plural-publics}

Human authorship does not mean one human principal or one corporation
defines the good. High-impact AI should remain accountable to plural
publics, law, affected communities, domain expertise, and independent
review.

The system should not resolve moral disagreement by presenting one
generated answer as the view of humanity.

\subsection{13.10 The external mirror and its
danger}\label{the-external-mirror-and-its-danger}

An earlier formulation imagined AI as an external self-deception
detector. this paper retains a weaker, safer possibility: AI can
function as a \textbf{justificatory mirror} by asking structured
questions and presenting counterevidence.

It should not diagnose hidden motives or certify moral purity. Such a
system could be used to pathologize dissent, manipulate vulnerable
users, or place institutional ideology behind a technical interface.

A defensible moral-assistance system should:

\begin{itemize}
\tightlist
\item
  show sources and uncertainty;
\item
  distinguish evidence from interpretation;
\item
  present multiple plausible objections;
\item
  disclose system and institutional constraints;
\item
  permit refusal;
\item
  preserve user privacy;
\item
  avoid claims of therapeutic or spiritual authority;
\item
  and provide appeal where outputs affect rights or access.
\end{itemize}

\section{14. Research program and principle
architecture}\label{research-program-and-principle-architecture}

\subsection{14.1 Empirical program}\label{empirical-program}

The framework generates research questions rather than one validation
test.

\subsubsection{Justificatory integrity}\label{justificatory-integrity}

\begin{itemize}
\tightlist
\item
  Do preregistered decision criteria reduce identity-sensitive standard
  shifting?
\item
  Does independent evidence preservation improve correction after
  institutional failure?
\item
  Which deliberative formats reduce motivated asymmetry without
  suppressing legitimate value disagreement?
\end{itemize}

\subsubsection{Power and participation}\label{power-and-participation}

\begin{itemize}
\tightlist
\item
  When does affected-community co-governance change which harms and
  alternatives are recognized?
\item
  Which participation structures create real authority rather than
  consultation theater?
\item
  How do retaliation risk and dependency alter apparent consent?
\end{itemize}

\subsubsection{Responsibility}\label{responsibility}

\begin{itemize}
\tightlist
\item
  How do people allocate blame when causal contribution, authority,
  benefit, and coercion diverge?
\item
  Which institutional designs prevent responsibility from concentrating
  on frontline workers while disappearing from leadership?
\item
  How do repair obligations track benefit and capacity independently of
  personal culpability?
\end{itemize}

\subsubsection{Moral injury}\label{moral-injury}

\begin{itemize}
\tightlist
\item
  Which forms of moral distress are primarily individual, relational,
  organizational, or political?
\item
  Which interventions restore trust without demanding premature
  reconciliation?
\item
  How do acknowledgment, material repair, and institutional change
  interact over time?
\end{itemize}

\subsubsection{AI ethics}\label{ai-ethics}

\begin{itemize}
\tightlist
\item
  Can AI-assisted counterargument improve reason quality without
  increasing conformity to system framing?
\item
  Which interface designs preserve meaningful refusal and human review?
\item
  How should uncertainty about artificial patiency alter experimentation
  and deployment?
\end{itemize}

\subsection{14.2 Falsifiability and abandonment
conditions}\label{falsifiability-and-abandonment-conditions}

The account should be revised or abandoned if evidence shows that:

\begin{itemize}
\tightlist
\item
  the proposed integrity diagnostics do not distinguish motivated
  asymmetry from ordinary error;
\item
  affected-party participation systematically worsens recognition and
  correction under conditions where stronger alternatives exist;
\item
  the responsibility profile omits variables necessary across cases or
  produces persistent misattribution;
\item
  practical correctability does not improve trust, error correction, or
  recurrence outcomes;
\item
  the moral-injury distinctions fail to support reliable research or
  intervention;
\item
  or the framework's governance rules are routinely used to silence
  rather than protect affected persons.
\end{itemize}

Normative principles are not falsified in exactly the same manner as
empirical hypotheses. Their coherence, implications, and public
defensibility can nevertheless be challenged by counterexample,
contradiction, exclusion, or intolerable consequence.

\subsection{14.3 Candidate Ethics
principles}\label{candidate-ethics-principles}

the account proposes the following family for later Foundations
consolidation. All remain provisional.

\subsubsection{CU-E-1 --- The Justificatory Integrity
Constraint}\label{cu-e-1-the-justificatory-integrity-constraint}

\textbf{{[}Normative principle{]}} An action, policy, or institution
lacks adequate moral justification when its defense requires systematic
distortion of relevant facts, reasons, affected parties, alternatives,
or standards in order to protect a dominant identity, advantage,
authority, or prior commitment. Integrity is necessary but not
sufficient for permissibility.

\subsubsection{CU-E-2 --- Agency, Patiency, and Degrees of
Responsibility}\label{cu-e-2-agency-patiency-and-degrees-of-responsibility}

\textbf{{[}Normative principle{]}} Moral patiency, moral agency,
responsibility, and sovereignty are distinct. Responsibility varies with
understanding, reason-responsiveness, control, authorship, coercion, and
capacity for repair; diminished agency can reduce blame while increasing
duties of protection.

\subsubsection{CU-E-3 --- Recognition and Affected-Being
Standing}\label{cu-e-3-recognition-and-affected-being-standing}

\textbf{{[}Normative principle{]}} Ethical evaluation must represent all
materially affected beings according to relevant welfare, agency,
vulnerability, relationship, history, and stake rather than treating
absence of institutional voice as absence of a claim.

\subsubsection{CU-E-4 --- Diachronic and Beneficiary
Responsibility}\label{cu-e-4-diachronic-and-beneficiary-responsibility}

\textbf{{[}Normative principle{]}} Moral evaluation includes delayed,
cumulative, intergenerational, and path-dependent consequences, together
with responsibilities created by prior action, inherited advantage,
ongoing benefit, and present capacity to prevent or repair harm.

\subsubsection{CU-E-5 --- Self-Deception as Moral
Failure}\label{cu-e-5-self-deception-as-moral-failure}

\textbf{{[}Normative principle{]}} Motivated distortion becomes morally
significant when it conceals avoidable harm, excludes affected beings,
corrupts standards, blocks correction, diffuses responsibility, or
protects power from accountability. Disagreement or error alone does not
establish self-deception.

\subsubsection{CU-E-6 --- Protection, Accountability, and
Repair}\label{cu-e-6-protection-accountability-and-repair}

\textbf{{[}Normative principle{]}} Responses to wrongdoing should
prioritize cessation and protection, truthful acknowledgment,
proportionate accountability, restitution or reparation, restoration of
standing, and prevention of recurrence. Punishment requires separate
justification and is neither intrinsically reparative nor categorically
excluded.

\subsubsection{CU-E-7 --- Moral Injury}\label{cu-e-7-moral-injury}

\textbf{{[}Bridge principle{]}} Moral injury names enduring
psychological, social, relational, existential, or spiritual suffering
associated with perpetrating, witnessing, failing to prevent, suffering,
or being betrayed by morally significant wrong. Its descriptive use
remains governed by the developing empirical literature.

\subsubsection{CU-E-8 --- Structural Moral Injury and Institutional
Betrayal}\label{cu-e-8-structural-moral-injury-and-institutional-betrayal}

\textbf{{[}Normative and empirical hypothesis{]}} Institutions can
produce or intensify moral injury when they impose morally compromising
roles, deny adequate authority or alternatives, conceal harm, betray
dependence, punish disclosure, or individualize structurally generated
distress.

\subsubsection{CU-E-9 --- Local--Global Compatibility and
Noncompensation}\label{cu-e-9-localglobal-compatibility-and-noncompensation}

\textbf{{[}Normative principle{]}} Claims of collective benefit are
inadequate when they depend on avoidable destruction of protected
standing, agency, or welfare, or on hidden transfer of severe,
nonconsensual, concentrated, or irreparable burdens. Sacrifice requires
necessity, proportionality, distributional analysis, and repair.

\subsubsection{CU-E-10 --- Political Legitimacy and Public
Justification}\label{cu-e-10-political-legitimacy-and-public-justification}

\textbf{{[}Normative principle{]}} Coercive or high-impact power is
legitimate only to the degree that it is publicly answerable to those
subject to it, preserves protected standing, remains contestable,
discloses evidence and uncertainty, and includes practical correction
and remedy.

\subsubsection{CU-E-11 --- Rights as Protected Conditions of Agency and
Patiency}\label{cu-e-11-rights-as-protected-conditions-of-agency-and-patiency}

\textbf{{[}Normative principle{]}} Rights protect bodily and
psychological integrity, welfare, liberty, capability, participation,
and standing against domination and unrestricted aggregate tradeoff.
Their content requires independent moral, legal, historical, and
democratic argument.

\subsubsection{CU-E-12 --- Democracy as a Conditional Corrective
Achievement}\label{cu-e-12-democracy-as-a-conditional-corrective-achievement}

\textbf{{[}Normative principle{]}} Democratic forms are valuable when
they distribute political agency, preserve opposition and minority
standing, constrain power, and enable peaceful correction. Majority rule
or actual consensus alone is insufficient.

\subsubsection{CU-E-13 --- Constraint on Political
Means}\label{cu-e-13-constraint-on-political-means}

\textbf{{[}Normative principle{]}} Political ends do not justify means
that destroy the truth conditions, protected standing, accountability,
or corrective institutions required to pursue, evaluate, and revise
those ends.

\subsubsection{CU-E-14 --- Responsibility--Authority
Alignment}\label{cu-e-14-responsibilityauthority-alignment}

\textbf{{[}Normative principle{]}} Responsibility should be assigned
with attention to causal contribution, knowledge, control, authority,
benefit, coercion, prevention capacity, and repair capacity.
Institutions should neither impose responsibility without adequate
authority nor permit authority without accountability.

\subsubsection{CU-E-15 --- Affected-Community
Participation}\label{cu-e-15-affected-community-participation}

\textbf{{[}Methodological and normative principle{]}} People governed,
classified, burdened, or represented by an ethical principle must have
substantive, accessible, and appropriately authoritative participation
in its development, application, and review. Participation supplements
rather than replaces evidence and argument. Promotion into Foundations
requires closure of every review gate assigned to the principle;
unresolved disagreement among required constituencies keeps the
principle provisional, with dissent and dispositions preserved in the
public record.

\subsubsection{CU-E-16 --- Moral Uncertainty, Precaution, and
Reversibility}\label{cu-e-16-moral-uncertainty-precaution-and-reversibility}

\textbf{{[}Normative principle{]}} Under credible risk of severe or
irreversible harm, uncertainty should alter action in proportion to
severity, plausibility, alternatives, reversibility, and the harms of
delay. An application must either protect reversibility, containment,
rollback, staged exposure, or equivalent option value through a declared
floor, or publicly justify a necessity-and-proportionality exception
that addresses safer alternatives and the harms of delay. Uncertainty
neither automatically permits exposure nor requires paralysis.

\subsubsection{CU-E-17 --- Moral Patiency Is Not
Sovereignty}\label{cu-e-17-moral-patiency-is-not-sovereignty}

\textbf{Status:} SUPERSEDED ALIAS; controlled by CU-G6.

\textbf{{[}Normative principle{]}} Evidence that a being can experience
welfare or suffering can ground protection without establishing
responsibility, authority, or legitimate control over others.

\subsubsection{CU-E-18 --- Compassionate Moral
Address}\label{cu-e-18-compassionate-moral-address}

\textbf{{[}Normative commitment{]}} Accountability should preserve the
standing of persons wherever protection allows, distinguish judgment
from humiliation, and seek understanding sufficient for prevention and
repair without converting compassion into excuse.

\subsection{14.4 Inclusion, validation, and consequential
authorization}\label{inclusion-validation-and-consequential-authorization}

Canonical inclusion means that an entry is the controlling CU
formulation of what the framework presently proposes. It does not mean
that philosophical, legal, empirical, clinical, or affected-community
review is complete, and it does not authorize the principle's
consequential use over people or institutions.

The complete v0.10 framework should be available for review in its
strongest internally disciplined form. Subsequent expert and
participatory criticism may narrow, split, reclassify, or reject a
principle through a versioned correction. Review therefore corrects the
complete proposal rather than preventing the proposal from being stated.

Consequential application has a higher threshold. Principles concerning
disability, reproductive ethics, punishment, migration, labor,
racialized or gendered harm, animal status, moral injury, or artificial
patiency require field-specific evidence, legal and institutional review
where applicable, explicit affected-being analysis, counterexample and
misuse testing, practical correction, and substantive participation by
people directly affected by the classifications and practices at issue.

Unresolved disagreement must remain public. Canonical numbering does not
make any constituency infallible, but neither may institutional power
convert publication status into permission to classify, burden, or
govern others without standing, safety, challenge, and remedy.

\section{15. What commitments this account
licenses}\label{what-commitments-this-account-licenses}

\subsection{15.1 Register declaration}\label{register-declaration}

The commitments in this section are not empirical discoveries, legal
verdicts, or deductions from physical law. They are official
commitment-register orientations constrained by the descriptive and
normative analysis developed above. They remain answerable to facts,
affected-party testimony, consequence, criticism, and revision.

\subsection{15.2 The descriptive and normative
ground}\label{the-descriptive-and-normative-ground}

Human moral judgment is vulnerable to motivated distortion. Power shapes
which facts and people become visible. Sincerity does not guarantee
justification. Aggregate benefit can conceal concentrated harm.
Responsibility is distributed across causal contribution, knowledge,
control, authority, benefit, and repair capacity. Institutions can
preserve order while blocking correction. Repair requires more than
apology or punishment.

The normative ground is inherited from the Mandala: beings with stakes,
experience, agency, vulnerability, and dependence inhabit relations in
which actions can protect, harm, dominate, recognize, or repair. The
ethical question is how to remain truthful and accountable within those
relations.

\subsection{15.3 Ethical truthfulness}\label{ethical-truthfulness}

Ethical practice is the willingness to let truth change who we are.

It requires more than factual accuracy. It asks whether our identities
depend upon another person's silence, whether our comfort depends upon
costs we do not see, whether loyalty has changed our evidential
standards, whether our abstractions erase the people they govern, and
whether our explanations have become instruments for avoiding
responsibility.

To live ethically is not to become free of bias, defensiveness, or
self-interest. It is to build a life and institutions in which those
forces do not possess final authority over what can be known, admitted,
challenged, or repaired.

\subsection{15.4 Truth before innocence}\label{truth-before-innocence}

The first impulse after moral criticism is often to establish that one
is not the kind of person who would do the thing alleged. That impulse
can make identity more important than reality.

This account licenses a different orientation:

\begin{quote}
Seek the truth about the action before seeking innocence for the self.
\end{quote}

A truthful response may discover fault, limited control, unintended
consequence, inherited structure, mixed motive, or genuine disagreement.
The point is not compulsory guilt. It is refusing to decide in advance
that self-preservation will govern the inquiry.

The same applies institutionally. A system should not begin from the
requirement that its policy, mission, profession, nation, party,
laboratory, or tradition remain innocent. It should begin from the
requirement that consequences and responsibilities become knowable.

\subsection{15.5 Compassion without
excuse}\label{compassion-without-excuse}

Compassion widens moral attention. It makes another person's
vulnerability, history, dependence, and suffering available to judgment.
It also makes it possible to see the fear, shame, and defensive
structure of a wrongdoer without reducing the wrong to those causes.

Compassion does not require suspended evidence, erased boundaries,
restored trust, reconciliation, or exemption from accountability. Its
ethical function is to prevent judgment from becoming humiliation and to
keep protection and repair oriented toward beings rather than toward
revenge or moral spectacle.

\subsection{15.6 Truth after wrongdoing}\label{truth-after-wrongdoing}

When truth reveals harm, the commitment is not merely to admit it.
Ethical truthfulness asks what must stop, what must be protected, whose
record must be corrected, what restitution is owed, whose standing must
be restored, and what structures must change.

Repair may include apology, punishment, compensation, care,
institutional reform, memorialization, boundary, separation, or refusal.
No single response is always sufficient. Forgiveness and reconciliation
may be valuable, impossible, inappropriate, or unsafe; they cannot be
demanded as proof of healing.

\subsection{15.7 CU-COM14 --- Love as the Ground and Aim of Conscious
Life}\label{cu-com14-love-as-the-ground-and-aim-of-conscious-life}

Love is not a substitute for justice. It is the commitment to let
another being's reality, vulnerability, freedom, and flourishing matter
in itself rather than only as an input to one's own project. In ethical
life, love widens standing and motivates care while truth and protected
conditions prevent care from becoming possession, paternalism, or
self-erasure.

\textbf{CU-T12 --- Love as the Relational Interior of Coherence} is an
interpretation. \textbf{CU-COM14 --- Love as the Ground and Aim of
Conscious Life} is the corresponding undertaking. Neither licenses
compulsory intimacy, forced forgiveness, reconciliation with abuse, or
sacrificial demands imposed by the powerful upon the vulnerable.

\subsection{15.8 CU-COM17 --- Redemptive Possibility Without
Justification of
Evil}\label{cu-com17-redemptive-possibility-without-justification-of-evil}

Ethics preserves the possibility that harm can be transformed into
solidarity, wisdom, creation, justice, care, or repair. That
transformation does not reach backward to make the harm necessary,
deserved, or good.

\begin{quote}
Evil may be redeemed; it is never retroactively justified.
\end{quote}

No victim owes redemption, meaning, testimony, forgiveness, or growth. A
redemptive response must remain answerable to protection, truth,
restitution, changed structures, and the authority of those who bore the
harm.

\subsection{15.9 CU-COM19 --- Enact Omega
here}\label{cu-com19-enact-omega-here}

At the ethical scale, the Immanent Omega is not mystical status or final
perfection. It is universal concern made concrete in a particular
relation: truth told without domination, power limited for the sake of
another's agency, suffering answered with care, and harm met with
repair.

A genuine act of love can be complete in kind while partial in extent.
It does not finish moral history or confer authority on the actor. The
test of Omega embodiment is whether conduct becomes more truthful,
humble, non-dominating, accountable, and capable of repair.

\subsection{15.10 CU-COM12 --- Ethical Truthfulness and
Accountability}\label{cu-com12-ethical-truthfulness-and-accountability}

\begin{quote}
Coherence Universalism commits to seeking truth about actions before
seeking innocence for agents or institutions. This includes truth about
motive, bias, identity protection, benefit, power, omitted alternatives,
displaced costs, affected-party testimony, responsibility, and what
repair requires.
\end{quote}

\textbf{Practical expression.} Apply symmetrical evidential standards;
disclose motive, benefit, power, uncertainty, and omitted alternatives;
hear affected-party testimony; preserve challengeable records; align
responsibility with authority and causal reach; and join acknowledgment
to practical repair.

\textbf{Boundary conditions.} CU-COM12 does not license mind reading,
compulsory confession, diagnosis of opponents, substitution of sincerity
for rights, or the use of self-deception language to silence
disagreement. Truthfulness is necessary but not sufficient for justice.

\textbf{Revision conditions.} The commitment must be narrowed or revised
where its concepts become clinically stigmatizing, coercive,
epistemically unreliable, or systematically easier for the powerful to
impose on others than upon themselves.

\subsection{15.11 A vow of moral truthfulness and
repair}\label{a-vow-of-moral-truthfulness-and-repair}

\begin{quote}
We commit to seek the truth about our actions before seeking innocence
for ourselves.\\
We commit to hear the people who bear consequences our models omit.\\
We commit to submit power to correction that power does not control.\\
We commit to distinguish explanation from excuse.\\
We commit to revise our self-understanding when reality contradicts
it.\\
We commit to protect standing, agency, and welfare where our power can
affect them.\\
And when truth reveals harm, we commit not merely to confess it, but to
repair what can still be repaired.
\end{quote}

\subsection{15.12 Boundaries of the
commitment}\label{boundaries-of-the-commitment}

This account does not sanctify accusation, demand endless
self-suspicion, or treat moral life as a purity project. It does not
require victims to educate wrongdoers, disclose trauma, forgive,
reconcile, or remain in relationship. It does not infer moral guilt from
distress, neurodivergence, cultural difference, or error. It does not
grant any institution the authority to define another person's inner
truth.

The commitment is to reality, not to confession as ritual. Its test is
whether persons and institutions become more truthful, less dominating,
more correctable, and more capable of protection and repair.

\section{16. Conclusion: truth before
innocence}\label{conclusion-truth-before-innocence}

Ethics is often presented as a search for the right rule, the correct
outcome, the proper virtue, or the legitimate procedure. Each of those
matters. But none can function well when agents and institutions control
in advance which facts, people, motives, costs, and alternatives are
permitted to appear.

Coherence Universalism therefore locates a central ethical task in the
preservation of truthful contact. Moral judgment must remain open to
realities that threaten identity, advantage, loyalty, and power. It must
represent affected beings, preserve protected conditions, disclose
uncertainty, examine distributions, review option sets, align
responsibility with authority, and remain practically correctable.

Justificatory integrity is necessary because corrupted reasoning cannot
reliably answer to those realities. It is insufficient because an honest
injustice remains an injustice. Ethics requires both truthful
justification and substantive regard for standing, agency, welfare,
rights, distribution, legitimacy, reversibility, and repair.

The deepest return of the series is now visible. Epistemology teaches us
not to lie about the universe. Ethics teaches us not to lie about
ourselves. The first commitment protects description from desire. The
second protects moral life from the identities and institutions that
would rather preserve innocence than face consequence.

Truth is not the end of ethics. It is the condition under which
responsibility and repair can begin.

\section{Appendix A. Ethics principles and application
governance}\label{appendix-a.-ethics-principles-and-application-governance}

\subsection{A.1 Domain authority and
status}\label{a.1-domain-authority-and-status}

CU-E-1 through CU-E-18 are the current Ethics-owned principle
formulations. They remain provisional normative, bridge, empirical, or
methodological principles as indicated in the technical registry. Their
presence here does not make them scientific laws or substitute for
philosophical, legal, empirical, and affected-community review.

\subsubsection{CU-E-1 --- The Justificatory Integrity
Constraint}\label{cu-e-1-the-justificatory-integrity-constraint-1}

An action, policy, or institution lacks adequate moral justification
when its defense requires systematic distortion of relevant facts,
reasons, affected parties, alternatives, or standards in order to
protect a dominant identity, advantage, authority, or prior commitment.
Integrity is necessary but not sufficient for permissibility.

\subsubsection{CU-E-2 --- Agency, Patiency, and Degrees of
Responsibility}\label{cu-e-2-agency-patiency-and-degrees-of-responsibility-1}

Moral patiency, moral agency, responsibility, and sovereignty are
distinct. Responsibility varies with understanding,
reason-responsiveness, control, authorship, coercion, and capacity for
repair; diminished agency can reduce blame while increasing duties of
protection.

\subsubsection{CU-E-3 --- Recognition and Affected-Being
Standing}\label{cu-e-3-recognition-and-affected-being-standing-1}

Ethical evaluation must represent all materially affected beings
according to relevant welfare, agency, vulnerability, relationship,
history, and stake rather than treating absence of institutional voice
as absence of a claim.

\subsubsection{CU-E-4 --- Diachronic and Beneficiary
Responsibility}\label{cu-e-4-diachronic-and-beneficiary-responsibility-1}

Moral evaluation includes delayed, cumulative, intergenerational, and
path-dependent consequences, together with responsibilities created by
prior action, inherited advantage, ongoing benefit, and present capacity
to prevent or repair harm.

\subsubsection{CU-E-5 --- Self-Deception as Moral
Failure}\label{cu-e-5-self-deception-as-moral-failure-1}

Motivated distortion becomes morally significant when it conceals
avoidable harm, excludes affected beings, corrupts standards, blocks
correction, diffuses responsibility, or protects power from
accountability. Disagreement or error alone does not establish
self-deception.

\subsubsection{CU-E-6 --- Protection, Accountability, and
Repair}\label{cu-e-6-protection-accountability-and-repair-1}

Responses to wrongdoing should prioritize cessation and protection,
truthful acknowledgment, proportionate accountability, restitution or
reparation, restoration of standing, and prevention of recurrence.
Punishment requires separate justification and is neither intrinsically
reparative nor categorically excluded.

\subsubsection{CU-E-7 --- Moral Injury}\label{cu-e-7-moral-injury-1}

Moral injury names enduring psychological, social, relational,
existential, or spiritual suffering associated with perpetrating,
witnessing, failing to prevent, suffering, or being betrayed by morally
significant wrong. Its descriptive use remains governed by the
developing empirical literature.

\subsubsection{CU-E-8 --- Structural Moral Injury and Institutional
Betrayal}\label{cu-e-8-structural-moral-injury-and-institutional-betrayal-1}

Institutions can produce or intensify moral injury when they impose
morally compromising roles, deny adequate authority or alternatives,
conceal harm, betray dependence, punish disclosure, or individualize
structurally generated distress.

\subsubsection{CU-E-9 --- Local--Global Compatibility and
Noncompensation}\label{cu-e-9-localglobal-compatibility-and-noncompensation-1}

Claims of collective benefit are inadequate when they depend on
avoidable destruction of protected standing, agency, or welfare, or on
hidden transfer of severe, nonconsensual, concentrated, or irreparable
burdens. Sacrifice requires necessity, proportionality, distributional
analysis, and repair.

\subsubsection{CU-E-10 --- Political Legitimacy and Public
Justification}\label{cu-e-10-political-legitimacy-and-public-justification-1}

Coercive or high-impact power is legitimate only to the degree that it
is publicly answerable to those subject to it, preserves protected
standing, remains contestable, discloses evidence and uncertainty, and
includes practical correction and remedy.

\subsubsection{CU-E-11 --- Rights as Protected Conditions of Agency and
Patiency}\label{cu-e-11-rights-as-protected-conditions-of-agency-and-patiency-1}

Rights protect bodily and psychological integrity, welfare, liberty,
capability, participation, and standing against domination and
unrestricted aggregate tradeoff. Their content requires independent
moral, legal, historical, and democratic argument.

\subsubsection{CU-E-12 --- Democracy as a Conditional Corrective
Achievement}\label{cu-e-12-democracy-as-a-conditional-corrective-achievement-1}

Democratic forms are valuable when they distribute political agency,
preserve opposition and minority standing, constrain power, and enable
peaceful correction. Majority rule or actual consensus alone is
insufficient.

\subsubsection{CU-E-13 --- Constraint on Political
Means}\label{cu-e-13-constraint-on-political-means-1}

Political ends do not justify means that destroy the truth conditions,
protected standing, accountability, or corrective institutions required
to pursue, evaluate, and revise those ends.

\subsubsection{CU-E-14 --- Responsibility--Authority
Alignment}\label{cu-e-14-responsibilityauthority-alignment-1}

Responsibility should be assigned with attention to causal contribution,
knowledge, control, authority, benefit, coercion, prevention capacity,
and repair capacity. Institutions should neither impose responsibility
without adequate authority nor permit authority without accountability.

\subsubsection{CU-E-15 --- Affected-Community
Participation}\label{cu-e-15-affected-community-participation-1}

People governed, classified, burdened, or represented by an ethical
principle must have substantive, accessible, and appropriately
authoritative participation in its development, application, and review.
Participation supplements rather than replaces evidence and argument.
Promotion requires closure of every assigned review gate; unresolved
conflict preserves provisional status and public dissent.

\subsubsection{CU-E-16 --- Moral Uncertainty, Precaution, and
Reversibility}\label{cu-e-16-moral-uncertainty-precaution-and-reversibility-1}

Under credible risk of severe or irreversible harm, uncertainty should
alter action in proportion to severity, plausibility, alternatives,
reversibility, and the harms of delay. The application must either
declare a context-appropriate reversibility, containment, rollback,
staged-exposure, or retained-option floor, or publicly justify a
necessity-and-proportionality exception. Uncertainty neither
automatically permits exposure nor requires paralysis.

\subsubsection{CU-E-17 --- Moral Patiency Is Not
Sovereignty}\label{cu-e-17-moral-patiency-is-not-sovereignty-1}

\textbf{Status:} SUPERSEDED ALIAS; controlled by CU-G6.

Evidence that a being can experience welfare or suffering can ground
protection without establishing responsibility, authority, or legitimate
control over others.

\subsubsection{CU-E-18 --- Compassionate Moral
Address}\label{cu-e-18-compassionate-moral-address-1}

Accountability should preserve the standing of persons wherever
protection allows, distinguish judgment from humiliation, and seek
understanding sufficient for prevention and repair without converting
compassion into excuse.

\subsubsection{CU-COM12 --- Ethical Truthfulness and
Accountability}\label{cu-com12-ethical-truthfulness-and-accountability-1}

Coherence Universalism commits to seeking truth about actions before
seeking innocence for agents or institutions. This includes truth about
motive, bias, identity protection, benefit, power, omitted alternatives,
displaced costs, affected-party testimony, responsibility, and what
repair requires.

\subsection{A.2 Twenty application-governance
rules}\label{a.2-twenty-application-governance-rules}

Any CU ethical application should satisfy the following rules.

\begin{enumerate}
\def\labelenumi{\arabic{enumi}.}
\tightlist
\item
  Name every materially affected group, including excluded and future
  parties.
\item
  Distinguish moral patiency, agency, responsibility, and authority.
\item
  State the descriptive facts and uncertainty separately from normative
  commitments.
\item
  Identify who defines the problem and candidate options, which
  alternatives were excluded, and whether affected parties can propose
  alternatives.
\item
  Map authority, ownership, sanctions, dependencies, and exit costs.
\item
  Report distributions, not only aggregates.
\item
  Declare protected floors and the argument for each.
\item
  Separate consent from compliance, dependency, and absence of
  resistance.
\item
  Treat justificatory integrity as necessary but not sufficient.
\item
  Do not infer self-deception from disagreement, error, identity, or
  distress alone.
\item
  Apply evidential and justificatory standards symmetrically to allies
  and opponents.
\item
  Distinguish causal contribution, answerability, blame, liability, and
  remedial duty.
\item
  Align responsibility with knowledge, control, authority, benefit, and
  repair capacity.
\item
  Preserve dissent, provenance, and independent evidence.
\item
  State the time horizon, tail-risk rule, reversibility assumptions, and
  whether reversibility functions as a floor, comparative consideration,
  or justified exception.
\item
  Do not treat a complaint channel as evidence of practical
  correctability without access, independence, remedy, and recurrence
  data.
\item
  Do not require forgiveness, reconciliation, trust, or renewed contact
  as conditions of repair.
\item
  Give people governed or classified by a principle substantive standing
  in its development and review.
\item
  Keep spiritual interpretation and contemplative commitment distinct
  from empirical and public-normative claims.
\item
  Treat every screening and ranking result as conditional on the
  presented option set; floor compliance does not certify option-set
  adequacy.
\end{enumerate}

\section{Appendix B. Major revisions from
v0.09}\label{appendix-b.-major-revisions-from-v0.09}

\begin{longtable}[]{@{}
  >{\raggedright\arraybackslash}p{(\columnwidth - 6\tabcolsep) * \real{0.2500}}
  >{\raggedright\arraybackslash}p{(\columnwidth - 6\tabcolsep) * \real{0.2500}}
  >{\raggedright\arraybackslash}p{(\columnwidth - 6\tabcolsep) * \real{0.2500}}
  >{\raggedright\arraybackslash}p{(\columnwidth - 6\tabcolsep) * \real{0.2500}}@{}}
\toprule\noalign{}
\begin{minipage}[b]{\linewidth}\raggedright
Topic
\end{minipage} & \begin{minipage}[b]{\linewidth}\raggedright
v0.09 formulation
\end{minipage} & \begin{minipage}[b]{\linewidth}\raggedright
v0.10 formulation
\end{minipage} & \begin{minipage}[b]{\linewidth}\raggedright
Disposition
\end{minipage} \\
\midrule\noalign{}
\endhead
\bottomrule\noalign{}
\endlastfoot
Normative center & One negative principle centered on self-deception &
Truth-centered plural constraint architecture & Reconstructed \\
Justificatory integrity & Nearly complete criterion of impermissibility
& Necessary but insufficient condition of defensible judgment &
Narrowed \\
Moral grounding & Extended transcendental defense inside Ethics &
Inherited from the Mandala and Foundations; Ethics studies truthfulness
within moral reality & Relocated \\
Coherence & Ethical value closely tied to system coherence & No
universal moral scalar; organization can be coherent and unjust &
Retired as universal metric \\
Consequences & Integrated into coherence-based analysis & Evaluated
directly with distribution, floors, uncertainty, and noncompensation &
Clarified \\
Rights & Derived through agency conditions & Explicit protected
conditions requiring independent normative and legal argument &
Reconstructed \\
Self-deception & Primary moral failure & Central but evidentially
demanding; disagreement and error do not suffice & Retained and
disciplined \\
Power & Important application context & Constitutive of framing,
evidence, options, voice, and remedy & Elevated \\
Responsibility & Primarily agent-centered & Multidimensional and
distributed across knowledge, control, authority, benefit, and repair
capacity & Expanded \\
Reversibility & Comparative consideration & Floor-or-exception
requirement under credible severe-harm risk & Added \\
Option sets & Candidate actions assumed & Screening is conditional on
presented alternatives; option-set adequacy requires independent review
& Added \\
Repair & Favored over punishment & Profile including protection, truth,
accountability, restitution, standing, and non-recurrence & Expanded \\
AI & External self-deception detector & Bounded justificatory mirror
embedded in sociotechnical institutions & Narrowed \\
Spiritual practice & Supporting moral formation & Commitment-register
practice bounded by evidence, non-humiliation, and substantive justice &
Retained with scope \\
\end{longtable}

\section{Appendix C. Technical supplement and
reproducibility}\label{appendix-c.-technical-supplement-and-reproducibility}

\textbf{Released evidence package.} The current supplement is
\href{../supplements/08_Ethics_Technical_Supplement_v0.10.md}{08\_Ethics\_Technical\_Supplement\_v0.10.md}.
Its companion evidence directory is
\href{../evidence/08_Ethics/README.md}{08\_Ethics}. The supplement and
evidence directory expose the derivations, deterministic constructions,
scripts, machine-readable results, source ledger, and audit materials
referenced in this manuscript. They do not elevate a toy model,
simulation, or reconstructed check into target-domain empirical
validation.

The accompanying technical supplement preserves:

\begin{itemize}
\tightlist
\item
  the complete formal propositions and non-entailment results;
\item
  the deterministic policy, responsibility, repair, correctability, and
  emergency-governance models;
\item
  the exhaustive scalarization lattice and tie convention;
\item
  the adversarial reversibility and option-set tests;
\item
  the complete claim and evidence ledger;
\item
  the full principle registry and promotion gates;
\item
  the selective literature and operationalization ledger;
\item
  the glossary;
\item
  the independent audit and bounded-repair record;
\item
  and the source-to-publication transformation map.
\end{itemize}

The constructions are stipulated demonstrations of logical distinctions.
They are not empirical moral scores, clinical instruments, legal tests,
or automated decision procedures.

\section{References}\label{references}

Anderson, E. S. (1999). What is the point of equality? \emph{Ethics,
109}(2), 287--337.

Anscombe, G. E. M. (1958). Modern moral philosophy. \emph{Philosophy,
33}(124), 1--19.

Arendt, H. (1963). \emph{Eichmann in Jerusalem: A report on the banality
of evil}. Viking.

Aristotle. (2000). \emph{Nicomachean ethics} (R. Crisp, Trans.).
Cambridge University Press. (Original work composed ca. fourth century
BCE)

Bandura, A. (1999). Moral disengagement in the perpetration of
inhumanities. \emph{Personality and Social Psychology Review, 3}(3),
193--209.

Bandura, A., Barbaranelli, C., Caprara, G. V., \& Pastorelli, C. (1996).
Mechanisms of moral disengagement in the exercise of moral agency.
\emph{Journal of Personality and Social Psychology, 71}(2), 364--374.

Batson, C. D., Thompson, E. R., Seuferling, G., Whitney, H., \&
Strongman, J. A. (1999). Moral hypocrisy: Appearing moral to oneself
without being so. \emph{Journal of Personality and Social Psychology,
77}(3), 525--537.

Bovens, M. (1998). \emph{The quest for responsibility: Accountability
and citizenship in complex organisations}. Cambridge University Press.

Darley, J. M., \& Latané, B. (1968). Bystander intervention in
emergencies: Diffusion of responsibility. \emph{Journal of Personality
and Social Psychology, 8}(4), 377--383.

Darwall, S. (2006). \emph{The second-person standpoint: Morality,
respect, and accountability}. Harvard University Press.

Ditto, P. H., Celniker, J. B., Siddiqi, S. S., Güngör, M., \& Relihan,
D. P. (2025). Partisan bias in political judgment. \emph{Annual Review
of Psychology, 76}, 717--740.

Fischer, J. M., \& Ravizza, M. (1998). \emph{Responsibility and control:
A theory of moral responsibility}. Cambridge University Press.

Forst, R. (2012). \emph{The right to justification: Elements of a
constructivist theory of justice} (J. Flynn, Trans.). Columbia
University Press.

Fricker, M. (2007). \emph{Epistemic injustice: Power and the ethics of
knowing}. Oxford University Press.

Gilligan, C. (1982). \emph{In a different voice}. Harvard University
Press.

Griffin, B. J., Purcell, N., Burkman, K., Litz, B. T., Bryan, C. J.,
Schmitz, M., Villierme, C., Walsh, J., \& Maguen, S. (2019). Moral
injury: An integrative review. \emph{Journal of Traumatic Stress,
32}(3), 350--362.

Habermas, J. (1990). \emph{Moral consciousness and communicative action}
(C. Lenhardt \& S. W. Nicholsen, Trans.). MIT Press.

Haidt, J. (2001). The emotional dog and its rational tail: A social
intuitionist approach to moral judgment. \emph{Psychological Review,
108}(4), 814--834.

Held, V. (2006). \emph{The ethics of care: Personal, political, and
global}. Oxford University Press.

Herman, J. L. (2023). \emph{Truth and repair: How trauma survivors
envision justice}. Basic Books.

Jameton, A. (1984). \emph{Nursing practice: The ethical issues}.
Prentice-Hall.

Kant, I. (1997). \emph{Groundwork of the metaphysics of morals} (M.
Gregor, Trans.). Cambridge University Press. (Original work published
1785)

Kelman, H. C., \& Hamilton, V. L. (1989). \emph{Crimes of obedience:
Toward a social psychology of authority and responsibility}. Yale
University Press.

Korsgaard, C. M. (1996). \emph{The sources of normativity}. Cambridge
University Press.

Korsgaard, C. M. (2018). \emph{Fellow creatures: Our obligations to the
other animals}. Oxford University Press.

Kunda, Z. (1990). The case for motivated reasoning. \emph{Psychological
Bulletin, 108}(3), 480--498.

Laozi. (1963). \emph{Tao Te Ching} (D. C. Lau, Trans.). Penguin Books.
(Original work composed ca. fourth century BCE)

Lewis, A., Lee, H. S., Zabelski, S., \& Shields, M. C. (2025).
Institutional betrayal in inpatient psychiatry: Effects on trust and
engagement with care. \emph{Psychiatric Services, 76}(3), 218--225.

Litz, B. T., Stein, N., Delaney, E., Lebowitz, L., Nash, W. P., Silva,
C., \& Maguen, S. (2009). Moral injury and moral repair in war veterans:
A preliminary model and intervention strategy. \emph{Clinical Psychology
Review, 29}(8), 695--706.

McAdams, M., Henninger, M. W., Bloeser, K., \& McCarron, K. K. (2025).
Institutional betrayal in military and veteran populations: A systematic
scoping review. \emph{Journal of the American Psychiatric Nurses
Association, 31}(1), 8--22.

Medina, J. (2013). \emph{The epistemology of resistance}. Oxford
University Press.

Mill, J. S. (2001). \emph{Utilitarianism} (2nd ed.). Hackett. (Original
work published 1863)

Molendijk, T. (2025). Addressing moral injury in practice: Suggestions
for organisational, political and societal interventions. \emph{European
Journal of Psychotraumatology, 16}(1), 2541485.

Monin, B., \& Miller, D. T. (2001). Moral credentials and the expression
of prejudice. \emph{Journal of Personality and Social Psychology,
81}(1), 33--43.

Moore, C., Detert, J. R., Treviño, L. K., Baker, V. L., \& Mayer, D. M.
(2012). Why employees do bad things: Moral disengagement and unethical
organizational behavior. \emph{Personnel Psychology, 65}(1), 1--48.

National Commission for the Protection of Human Subjects of Biomedical
and Behavioral Research. (1979). \emph{The Belmont report: Ethical
principles and guidelines for the protection of human subjects of
research}. U.S. Department of Health and Human Services.

Nussbaum, M. C. (2006). \emph{Frontiers of justice: Disability,
nationality, species membership}. Belknap Press.

O'Neill, O. (2002). \emph{Autonomy and trust in bioethics}. Cambridge
University Press.

Pettit, P. (1997). \emph{Republicanism: A theory of freedom and
government}. Oxford University Press.

Rader, G. K. D. (2026). \emph{Coherence Universalism:
Physics---Coherence, Accessibility, and the Ordered Emergence of
Physical Description} (Version 0.10). Heaven≡Earth Foundation.

Rader, G. K. D. (2026). \emph{Coherence Universalism: Metaphysics and
Epistemology---Coherence Under Constraint, Explanatory Order, and
Warrant Across Registers} (Version 0.10). Heaven≡Earth Foundation.

Rader, G. K. D. (2026). \emph{Coherence Universalism:
Foundations---Principles, Domain Gateways, Commitments, and Constraints
on Responsible Application} (Version 0.10). Heaven≡Earth Foundation.

Rader, G. K. D. (2026). \emph{Coherence Universalism: Biology---Life as
Active Self-Maintenance: Viability, Individuation, and Multiscale
Regulation} (Version 0.10). Heaven≡Earth Foundation.

Rader, G. K. D. (2026). \emph{Coherence Universalism:
Consciousness---Perspective, Temporal Depth, Valence, and the
Architecture of Experience} (Version 0.10). Heaven≡Earth Foundation.

Rader, G. K. D. (2026). \emph{Coherence Universalism:
Psychology---Meaning, Regulation, and Repair Across Time} (Version
0.10). Heaven≡Earth Foundation.

Rader, G. K. D. (2026). \emph{Coherence Universalism: Social
Dynamics---Coordination, Power, Meaning, and Repair in Plural Societies}
(Version 0.10). Heaven≡Earth Foundation.

Rader, G. K. D. (2026). \emph{Denying Knowledge to Make Room for
Commitment---Truth, Commitment, and Nonmonotonic Warrant} (Version
0.10). Heaven≡Earth Foundation.

Rawls, J. (1971). \emph{A theory of justice}. Harvard University Press.

Rawls, J. (1993). \emph{Political liberalism}. Columbia University
Press.

Scanlon, T. M. (1998). \emph{What we owe to each other}. Belknap Press.

Sen, A. (2009). \emph{The idea of justice}. Harvard University Press.

Shay, J. (1994). \emph{Achilles in Vietnam: Combat trauma and the
undoing of character}. Scribner.

Silver, I., \& Berman, J. Z. (2024). What drives disagreement about
moral hypocrisy? Perceived comparability and how people exploit it to
criticize enemies and defend allies. \emph{Cognition, 247}, 105773.

Smith, C. P., \& Freyd, J. J. (2014). Institutional betrayal.
\emph{American Psychologist, 69}(6), 575--587.

Strawson, P. F. (1962). Freedom and resentment. \emph{Proceedings of the
British Academy, 48}, 1--25.

The Bhagavad Gita. (2007). (E. Easwaran, Trans.). Nilgiri Press.

Thompson, D. F. (1980). Moral responsibility of public officials: The
problem of many hands. \emph{American Political Science Review, 74}(4),
905--916.

Tronto, J. C. (1993). \emph{Moral boundaries: A political argument for
an ethic of care}. Routledge.

United Nations. (2005). \emph{Basic Principles and Guidelines on the
Right to a Remedy and Reparation for Victims of Gross Violations of
International Human Rights Law and Serious Violations of International
Humanitarian Law} (A/RES/60/147).

United Nations Committee on the Rights of Persons with Disabilities.
(2018). \emph{General comment No.~7 on the participation of persons with
disabilities, including children with disabilities, through their
representative organizations, in the implementation and monitoring of
the Convention} (CRPD/C/GC/7).

United Nations Office on Drugs and Crime. (2020). \emph{Handbook on
restorative justice programmes} (2nd ed.). United Nations.

UNESCO. (2005). \emph{Universal Declaration on Bioethics and Human
Rights}.

VanderWeele, T. J., Wortham, J. S., Carey, L. B., Case, B. W., Cowden,
R. G., Duffee, C., Jackson-Meyer, K., Lu, F., Mattson, S. A., Padgett,
R. N., Peteet, J. R., Rutledge, J., Symons, X., \& Koenig, H. G. (2025).
Moral trauma, moral distress, moral injury, and moral injury disorder:
Definitions and assessments. \emph{Frontiers in Psychology, 16},
1422441.

Walker, M. U. (2006). \emph{Moral repair: Reconstructing moral relations
after wrongdoing}. Cambridge University Press.

Williamson, V., Kothari, R., Bonson, A., Campbell, G., Greenberg, N.,
Murphy, D., \& Lamb, D. (2025). Moral injury prevention and
intervention. \emph{European Journal of Psychotraumatology, 16}(1),
2567721.

World Health Organization Regional Office for Europe. (2025).
\emph{Mental health of nurses and doctors survey in the European Union,
Iceland and Norway}. WHO.

Xue, K., Shang, J., Yang, C., Pan, L., Shi, H., \& Zeng, Y. (2025).
Nurses' moral distress and ethical climate: A systematic review and
meta-analysis. \emph{Nursing Ethics, 32}(7), 1981--1997.

Young, I. M. (2011). \emph{Responsibility for justice}. Oxford
University Press.

Zehr, H. (2002). \emph{The little book of restorative justice}. Good
Books.

\end{document}
